% !TEX program = lualatex
% ---------------------------------------------------------------------------------------------
%  Readback note 27: A Taylor–Wiles prime at a time: ties of tame inertia to Frobenius and the denominator of the projection
%
%  Build:  python3 tools/rb.py note <slug>
%
%  One unit of a Lean artifact, rewritten as ordinary mathematics. The prose is hand-written. Lean names,
%  line ranges and counts come from data.tex; quoted Lean statements are copied from the artifact by the build.
%
%  Rules of the text: standard terminology; every symbol is explained where it first appears and listed again
%  in the notation at the end; no letter has two meanings; a proof is complete or says what it takes for granted;
%  Lean names appear only in the LEAN lines and in the concordance.
%
%  The signature rule (the build enforces it): a function, relation or predicate that the artifact defines is
%  written only through its notation macro, which the build makes from the spec (`notation`) with one slot for
%  each explicit argument of the Lean definition. No argument is fixed and dropped, also not "for brevity" and
%  not for one section. Every argument and every hypothesis of a stated declaration stands inside its
%  statement, or in a Setting that the statement cites by number; only implicit and instance arguments (the
%  ambient objects) may stand in the conventions.
% ---------------------------------------------------------------------------------------------
\documentclass[11pt,a4paper]{article}
\ifdefined\pdfvariable\pdfvariable suppressoptionalinfo \numexpr 1+2+4+8\relax\fi
\usepackage[a4paper,hmargin=35mm,top=22mm,bottom=27mm,footskip=12mm,nohead]{geometry}

\usepackage{fontspec}
\directlua{
  luaotfload.add_fallback("gfallback", {"STIXTwoMath-Regular.otf:mode=node;", "DejaVuSans.ttf:mode=node;"})
  luaotfload.add_fallback("gmonofallback", {"DejaVuSansMono.ttf:mode=node;", "STIXTwoMath-Regular.otf:mode=node;"})
}
\setmainfont{NewCM10}[Extension=.otf,UprightFont=*-Book,ItalicFont=*-BookItalic,
  BoldFont=*-Bold,BoldItalicFont=*-BoldItalic,SmallCapsFont=*-Book,SmallCapsFeatures={Letters=SmallCaps},
  RawFeature={fallback=gfallback}]
\setsansfont{NewCMSans10}[Extension=.otf,UprightFont=*-Book,ItalicFont=*-BookOblique,
  BoldFont=*-Bold,BoldItalicFont=*-BoldOblique,RawFeature={fallback=gfallback}]
\setmonofont{JetBrainsMono}[Extension=.otf,UprightFont=*-Regular,ItalicFont=*-Italic,
  BoldFont=*-SemiBold,BoldItalicFont=*-SemiBoldItalic,Scale=MatchLowercase,
  RawFeature={-calt;-liga;fallback=gmonofallback}]
\usepackage{amsmath}
\usepackage{amsthm}
\usepackage[math-style=TeX]{unicode-math}
\setmathfont{NewCMMath-Book.otf}
\emergencystretch=2.5em\relax

\usepackage{xcolor}
\usepackage{tikz}
\usetikzlibrary{calc,arrows.meta,positioning,backgrounds}
\usepackage[protrusion=true,expansion=true]{microtype}
\usepackage{enumitem}
\usepackage{needspace}
\usepackage{booktabs}
\usepackage{array}
\usepackage{longtable}
\usepackage{fancyhdr}
\usepackage{lastpage}
\usepackage{multicol}
\usepackage[hidelinks,pdfusetitle]{hyperref}

\definecolor{ink}{HTML}{0F172A}
\definecolor{inksoft}{HTML}{475569}
\definecolor{inkfaint}{HTML}{94A3B8}
\definecolor{rulegrey}{HTML}{CBD5E1}
\definecolor{arith}{HTML}{C2410C}
\definecolor{est}{HTML}{7C3AED}
\definecolor{ana}{HTML}{0891B2}
\definecolor{ok}{HTML}{047857}
\definecolor{leancol}{HTML}{3730A3}

\title{A Taylor–Wiles prime at a time: ties of tame inertia to Frobenius and the denominator of the projection}
\author{grh-compact}
\date{}

% ------------------------------------------------------------------------------------- macros
\newcommand{\leanfont}{\ttfamily\fontsize{8.6}{10.5}\selectfont\color{leancol}}
\newcommand{\lean}[1]{\csname LN@#1\endcsname}
\newcommand{\fillin}[1]{{\color{arith}\itshape[#1]}}  % a place of the template that is not written yet; the build reports every one that is left
\newcommand{\leanquote}[1]{\csname LQ@#1\endcsname}
% generated by rb.py note; do not edit
\newcommand{\ConcDecls}{10}
\newcommand{\Pin}{72fead10}
\newcommand{\PinText}{that of commit {\leanfont 72fead10}}
\newcommand{\ArtifactName}{grh-compact}
\expandafter\newcommand\csname LN@lineStab\endcsname{\href{https://github.com/lsernau/grh-lean/blob/72fead102f663d5c7a407964a13dc0cf9d977207/grh-compact/src/New/Patching/OnePrime.lean\#L15-L50}{{\leanfont exists\_line\_stabilizer\_log}}}
\expandafter\newcommand\csname LN@weightedSum\endcsname{\href{https://github.com/lsernau/grh-lean/blob/72fead102f663d5c7a407964a13dc0cf9d977207/grh-compact/src/New/Patching/OnePrime.lean\#L61-L61}{{\leanfont weightedSum}}}
\expandafter\newcommand\csname LN@mulWeightedSum\endcsname{\href{https://github.com/lsernau/grh-lean/blob/72fead102f663d5c7a407964a13dc0cf9d977207/grh-compact/src/New/Patching/OnePrime.lean\#L64-L65}{{\leanfont mul\_weightedSum}}}
\expandafter\newcommand\csname LN@mapEqWeights\endcsname{\href{https://github.com/lsernau/grh-lean/blob/72fead102f663d5c7a407964a13dc0cf9d977207/grh-compact/src/New/Patching/OnePrime.lean\#L72-L74}{{\leanfont map\_eq\_of\_weights\_modEq}}}
\expandafter\newcommand\csname LN@firstOrder\endcsname{\href{https://github.com/lsernau/grh-lean/blob/72fead102f663d5c7a407964a13dc0cf9d977207/grh-compact/src/New/Patching/OnePrime.lean\#L81-L118}{{\leanfont one\_prime\_first\_order}}}
\expandafter\newcommand\csname LQ@firstOrder\endcsname{\leanquoteline{0}{lemma one\_prime\_first\_order (a : H →* Multiplicative (ZMod m)) (B : \_ →+ R) (hB : ∀ f (p : H),}\leanquoteline{2}{B f * ρ p = B (f <| · - (a p).toAdd)) (𝒜 : Subring R) (hρ : ∀ g, ρ g ∈ 𝒜) (h𝒜 : ∀ f, B f ∈ 𝒜) \{A\}}\leanquoteline{2}{\{ac : A → R\} (hc : ∀ \{x\}, x ∈ 𝒜 → ∀ t, Commute (ac t) x) \{Ψ L γ\} \{n k : ℕ\} (hΨ : ∑ p : H,}\leanquoteline{2}{ρ ↑p⁻¹ * Ψ * ρ p = n) (hX : ∀ p ∈ H, ρ p * B 1 = B 1) (h : L * B 1 = 0) (hL : ∀ g, L * ρ g =}\leanquoteline{2}{L) (tf : IsSMulRegular R (Fintype.card Γ)) (hn : k * (a.transfer γ).toAdd = k) : ∃ F ∈ 𝒜, (∀ g,}\leanquoteline{2}{∃ Z, σ (n • (F * (ρ g - 1))) = σ ((∑ g, ρ g) * Z) ∧ L * Z = 0) ∧ ∀ \{t t' r\}, (∀ f : ZMod m → ℤ,}\leanquoteline{2}{B ((f <| · - k) - f) * ac t = weightedSum (ρ.comp H.subtype) a ((f <| · - k) - f) * ac t') →}\leanquoteline{2}{(∀ f, weightedSum ρ a.transfer f * ac t' = (ρ γ - 1) * weightedSum ρ a.transfer f * r) →}\leanquoteline{2}{(2 * k) • (σ F * σ (ac t)) = (2 * k) • (σ (∑ g, ρ g) * σ r)}}
\expandafter\newcommand\csname LN@weightedSumExt\endcsname{\href{https://github.com/lsernau/grh-lean/blob/72fead102f663d5c7a407964a13dc0cf9d977207/grh-compact/src/New/Patching/OnePrimeTie.lean\#L44-L59}{{\leanfont BaseLevel.\allowbreak{}weightedSum\_ext}}}
\expandafter\newcommand\csname LN@charpolyMem\endcsname{\href{https://github.com/lsernau/grh-lean/blob/72fead102f663d5c7a407964a13dc0cf9d977207/grh-compact/src/New/Patching/OnePrimeTie.lean\#L66-L77}{{\leanfont BaseLevel.\allowbreak{}charpoly\_mem\_heckeAlgebra}}}
\expandafter\newcommand\csname LN@traceElem\endcsname{\href{https://github.com/lsernau/grh-lean/blob/72fead102f663d5c7a407964a13dc0cf9d977207/grh-compact/src/New/Patching/OnePrimeTie.lean\#L81-L85}{{\leanfont BaseLevel.\allowbreak{}traceElem}}}
\expandafter\newcommand\csname LN@existsTie\endcsname{\href{https://github.com/lsernau/grh-lean/blob/72fead102f663d5c7a407964a13dc0cf9d977207/grh-compact/src/New/Patching/OnePrimeTie.lean\#L96-L178}{{\leanfont BaseLevel.\allowbreak{}exists\_tie}}}
\expandafter\newcommand\csname LN@mainThm\endcsname{\href{https://github.com/lsernau/grh-lean/blob/72fead102f663d5c7a407964a13dc0cf9d977207/grh-compact/src/New/Patching/OnePrimeTie.lean\#L190-L304}{{\leanfont BaseLevel.\allowbreak{}exists\_diagOp\_mem\_ellIntegralForms}}}
\expandafter\newcommand\csname LQ@mainThm\endcsname{\leanquoteline{0}{lemma BaseLevel.exists\_diagOp\_mem\_ellIntegralForms (b : D.Eigensystem Q.w → PadicAlgCl ℓ) (p :}\leanquoteline{2}{EllAdicIntegers[X]) (𝓧 : Set (D.Eigensystem Q.w)) (π : 𝓧) e \{t\} (ht : t ≠ 0) (he :}\leanquoteline{2}{∀ j (x : D.Eigensystem Q.w), x.1 ∈ D.Eigensystem ![] → aeval (b x) p ≠ 0 → (∀ u : D.heckeAlgebra}\leanquoteline{2}{Q.w, u.1 x = 0 → ‖u.1 π‖ < 1) → HasRootString (x.1.heckePolyAt (Q.w j)) ((e j).map (algebraMap}\leanquoteline{2}{EllAdicIntegers ℂ)) (Q.w j).q (algebraMap EllAdicIntegers ℂ (a (Q.frobData j).frob))⁻¹ t) (hν :}\leanquoteline{2}{(2 * (card n)! : ℕ) * (p.coeff 0 * c \textasciicircum{} 2) \textasciicircum{} (card J + 1) ∣ ℓ \textasciicircum{} ν) \{m γ\} (rib : ∀ S 𝒥,}\leanquoteline{2}{IsTraceIdeal (S := S) (Y := D.Eigensystem Q.w) 𝒥 (N ·.1 · |>.1) (ℓ \textasciicircum{} ν : EllAdicIntegers)}\leanquoteline{2}{(Q.frobData · |>.isPlaceOver.tameGenerator) (Q.frobData · |>.frob) (·) e → ∃ d,}\leanquoteline{2}{IsRibetElement 𝒥 b π d ∧ ellVal (m * d) ≤ γ) : ∃ s, ellVal (m * s π.1) ≤}\leanquoteline{2}{γ ∧ ∀ l ∈ D.formsOfLevel ![], ∀ d : ℕ, d • l ∈ ellIntegralForms → ∀ y ∈ D.formsOfLevel ![],}\leanquoteline{2}{(∀ s, (∀ x ∈ 𝓧, s x = 0) → diagOp s y = 0) → (∀ s, (∀ x, aeval (b x) p = 0 → s x = 0) →}\leanquoteline{2}{diagOp s l = diagOp s y) → d • c • diagOp s y ∈ ellIntegralForms}}
\newcommand{\ConcordanceRows}{\crow{lineStab}{\href{https://github.com/lsernau/grh-lean/blob/72fead102f663d5c7a407964a13dc0cf9d977207/grh-compact/src/New/Patching/OnePrime.lean\#L15-L50}{{\leanfont exists\_line\_stabilizer\_log}}}{OnePrime.lean}{15--50}
\crow{weightedSum}{\href{https://github.com/lsernau/grh-lean/blob/72fead102f663d5c7a407964a13dc0cf9d977207/grh-compact/src/New/Patching/OnePrime.lean\#L61-L61}{{\leanfont weightedSum}}}{OnePrime.lean}{61--61}
\crow{mulWeightedSum}{\href{https://github.com/lsernau/grh-lean/blob/72fead102f663d5c7a407964a13dc0cf9d977207/grh-compact/src/New/Patching/OnePrime.lean\#L64-L65}{{\leanfont mul\_weightedSum}}}{OnePrime.lean}{64--65}
\crow{mapEqWeights}{\href{https://github.com/lsernau/grh-lean/blob/72fead102f663d5c7a407964a13dc0cf9d977207/grh-compact/src/New/Patching/OnePrime.lean\#L72-L74}{{\leanfont map\_eq\_of\_weights\_modEq}}}{OnePrime.lean}{72--74}
\crow{firstOrder}{\href{https://github.com/lsernau/grh-lean/blob/72fead102f663d5c7a407964a13dc0cf9d977207/grh-compact/src/New/Patching/OnePrime.lean\#L81-L118}{{\leanfont one\_prime\_first\_order}}}{OnePrime.lean}{81--118}
\crow{weightedSumExt}{\href{https://github.com/lsernau/grh-lean/blob/72fead102f663d5c7a407964a13dc0cf9d977207/grh-compact/src/New/Patching/OnePrimeTie.lean\#L44-L59}{{\leanfont BaseLevel.\allowbreak{}weightedSum\_ext}}}{OnePrimeTie.lean}{44--59}
\crow{charpolyMem}{\href{https://github.com/lsernau/grh-lean/blob/72fead102f663d5c7a407964a13dc0cf9d977207/grh-compact/src/New/Patching/OnePrimeTie.lean\#L66-L77}{{\leanfont BaseLevel.\allowbreak{}charpoly\_mem\_heckeAlgebra}}}{OnePrimeTie.lean}{66--77}
\crow{traceElem}{\href{https://github.com/lsernau/grh-lean/blob/72fead102f663d5c7a407964a13dc0cf9d977207/grh-compact/src/New/Patching/OnePrimeTie.lean\#L81-L85}{{\leanfont BaseLevel.\allowbreak{}traceElem}}}{OnePrimeTie.lean}{81--85}
\crow{existsTie}{\href{https://github.com/lsernau/grh-lean/blob/72fead102f663d5c7a407964a13dc0cf9d977207/grh-compact/src/New/Patching/OnePrimeTie.lean\#L96-L178}{{\leanfont BaseLevel.\allowbreak{}exists\_tie}}}{OnePrimeTie.lean}{96--178}
\crow{mainThm}{\href{https://github.com/lsernau/grh-lean/blob/72fead102f663d5c7a407964a13dc0cf9d977207/grh-compact/src/New/Patching/OnePrimeTie.lean\#L190-L304}{{\leanfont BaseLevel.\allowbreak{}exists\_diagOp\_mem\_ellIntegralForms}}}{OnePrimeTie.lean}{190--304}}
% the notations of the defined functions: one slot for each explicit argument of the Lean declaration
\newcommand{\weightedSum}[3]{\mathsf{W}_{#1,#2}(#3)}
\newcommand{\traceElem}[4]{\mathsf{t}_{#1}(#2;\,#3,#4)}
\newcommand{\UnitMembers}{10}
\newcommand{\UnitInNote}{10}
\newcommand{\UnitProofs}{6}
\newcommand{\UnitSteps}{122}
\newcommand{\UnitAllSteps}{144}


\linespread{1.08}                       % 11 pt on 14.7 pt
\setlength{\parindent}{0pt}
\setlength{\parskip}{5.5pt plus 1.5pt minus 0.5pt}
\setlength{\abovedisplayskip}{7.5pt plus 2pt minus 3pt}
\setlength{\belowdisplayskip}{7.5pt plus 2pt minus 3pt}
\setlength{\abovedisplayshortskip}{3pt plus 2pt}
\setlength{\belowdisplayshortskip}{6pt plus 2pt minus 2pt}
\emergencystretch=1.5em
\clubpenalty=10000 \widowpenalty=10000 \displaywidowpenalty=10000
\relpenalty=7000 \binoppenalty=9000
\raggedbottom
\pagestyle{fancy}\fancyhf{}\renewcommand{\headrulewidth}{0pt}
\fancyfoot[C]{\sffamily\fontsize{7.6}{9}\selectfont\color{inkfaint}\thepage\,/\,\pageref*{LastPage}}
\setlist{topsep=2pt,partopsep=0pt,parsep=0pt,itemsep=4pt,leftmargin=2.1em,labelsep=0.55em}

\newcommand{\Q}{\mathbb{Q}}
\newcommand{\Z}{\mathbb{Z}}
\newcommand{\F}{\mathbb{F}}
\newcommand{\len}{\operatorname{len}}
\newcommand{\cok}{\operatorname{cok}}
\newcommand{\End}{\operatorname{End}}
\newcommand{\Hom}{\operatorname{Hom}}
\newcommand{\Gal}{\operatorname{Gal}}
\newcommand{\diag}{\operatorname{diag}}
\newcommand{\kw}[1]{\emph{#1}}
\newcommand{\m}[1]{\mbox{$#1$}}        % a short formula that must not be broken across lines
\newcommand{\head}[1]{\par\addvspace{17pt plus 4pt minus 2pt}\Needspace{6\baselineskip}%
  \noindent{\fontsize{13.2}{16}\selectfont\bfseries\color{ink}#1}\par\nopagebreak\vspace{-1pt}}
\newcommand{\sect}[2]{\refstepcounter{section}\label{#2}\par\addvspace{19pt plus 4pt minus 2pt}\Needspace{7\baselineskip}%
  \noindent\llap{{\fontsize{13.2}{16}\selectfont\bfseries\color{inkfaint}\thesection}\hspace{9pt}}%
  {\fontsize{13.2}{16}\selectfont\bfseries\color{ink}#1}\par\nopagebreak\vspace{-1pt}}
\newcommand{\subhead}[1]{\par\addvspace{9pt plus 2pt}\Needspace{4\baselineskip}\noindent\textbf{#1}\enspace\ignorespaces}
% where a statement lives in the Lean archive
\newcommand{\leantag}{{\sffamily\bfseries\fontsize{6.3}{8}\selectfont\color{leancol}\addfontfeatures{LetterSpace=7}LEAN}}
\newcommand{\sep}{\unskip\nobreak\hspace{5pt}{\color{inkfaint}\textperiodcentered}\hspace{5pt plus 2pt}\ignorespaces}
\newcommand{\leanline}[1]{\par\nopagebreak\vspace{-2.5pt}{\fontsize{8.9}{12.4}\selectfont\color{inksoft}\raggedright
  \noindent\llap{\leantag\hspace{9pt}}#1\par}\vspace{1pt}}
% a displayed statement in a tinted box: #1 colour, #2 text
\newcommand{\statement}[2]{\par\addvspace{7pt}\begingroup\setlength{\fboxsep}{8.5pt}\setlength{\fboxrule}{0.6pt}%
  \noindent\fcolorbox{#1!30!white}{#1!3!white}{\parbox{\dimexpr\linewidth-2\fboxsep-2\fboxrule\relax}{#2}}%
  \endgroup\par\addvspace{7pt}}
\newcommand{\nt}[2]{\par\noindent\hangindent=1.1em\hangafter=1 #1\enspace{\color{inksoft}#2}\par\vspace{1pt}}
\newcommand{\legdot}[1]{\tikz[baseline=-0.62ex]\fill[#1] (0,0) circle (2.5pt);\hspace{3pt}}
\newcommand{\crow}[4]{\csname CS@#1\endcsname & #2 & \textcolor{inksoft}{#3} & \textcolor{inksoft}{#4}\\}
\newcommand{\leanquoteline}[2]{\par\noindent\hangindent=\dimexpr#1\fontdimen2\font+2.2em\relax\hangafter=1
  \hspace*{#1\fontdimen2\font}#2}
\newcommand{\quoted}[1]{\par\vspace{3pt}{\ttfamily\fontsize{7.5}{10.4}\selectfont\color{leancol}\raggedright
  \setlength{\parskip}{0pt}\leanquote{#1}\par}\vspace{4pt}}
\newcommand{\bx}[3]{{\sffamily\bfseries\fontsize{7.3}{9.2}\selectfont\color{#1!82!black}#2}\\[1.5pt]#3}
\newcommand{\cs}[2]{\expandafter\newcommand\csname CS@#1\endcsname{#2}}

% theorems
\newtheoremstyle{plainx}{7pt plus 2pt}{7pt plus 2pt}{\itshape}{}{\bfseries}{.}{0.5em}{}
\newtheoremstyle{uprightx}{7pt plus 2pt}{7pt plus 2pt}{}{}{\bfseries}{.}{0.5em}{}
\theoremstyle{plainx}
\newtheorem{theorem}{Theorem}[section]
\newtheorem{lemma}[theorem]{Lemma}
\newtheorem{proposition}[theorem]{Proposition}
\newtheorem{corollary}[theorem]{Corollary}
\theoremstyle{uprightx}
\newtheorem{definition}[theorem]{Definition}
\newtheorem{setting}[theorem]{Setting}          % hypotheses shared by several statements; each of them cites it by number
\newtheorem{remark}[theorem]{Remark}
\newtheorem{example}[theorem]{Example}
\renewcommand{\qedsymbol}{{\setlength{\fboxsep}{0pt}\setlength{\fboxrule}{0.5pt}\fbox{\rule{0pt}{0.46em}\hspace{0.46em}}}}
% what holds throughout the note: the ambient objects only (implicit and instance arguments), never a hypothesis
\newenvironment{conventions}{\par\addvspace{5pt}\noindent\textbf{Conventions.}\enspace\ignorespaces}{\par\addvspace{3pt}}
\makeatletter
\renewenvironment{proof}[1][Proof]{\par\pushQED{\qed}\normalfont\topsep0pt\partopsep0pt
  \trivlist\item[\hskip\labelsep\itshape #1.]\ignorespaces}{\popQED\endtrivlist\@endpefalse}
\makeatother

% where each declaration of the concordance is treated: \cs{key}{place}, the place naming a \label or a lettered item
\cs{lineStab}{Lemma~\ref{l:line}}
\cs{weightedSum}{Definition~\ref{d:wsum}}
\cs{mulWeightedSum}{Lemma~\ref{l:shift}}
\cs{mapEqWeights}{Lemma~\ref{l:modeq}}
\cs{firstOrder}{Proposition~\ref{p:first}}
\cs{weightedSumExt}{Lemma~\ref{l:test}}
\cs{charpolyMem}{Lemma~\ref{l:charpoly}}
\cs{traceElem}{Definition~\ref{d:trace}}
\cs{existsTie}{Proposition~\ref{p:tie}}
\cs{mainThm}{Theorem A}
\newcommand{\CC}{\mathbb{C}}
\newcommand{\ZZ}{\mathbb{Z}}
\newcommand{\NN}{\mathbb{N}}
\newcommand{\Zl}{\bar{\mathbb{Z}}_\ell}
\newcommand{\Ql}{\bar{\mathbb{Q}}_\ell}
\newcommand{\Qbar}{\bar{\mathbb{Q}}}
\newcommand{\GE}{G_E}
\newcommand{\Gf}{G_{\mathrm{f}}}
\newcommand{\Fw}[1]{\mathbb{F}_{#1}}
\newcommand{\Ew}[1]{E_{#1}}
\newcommand{\Ow}[1]{\mathcal{O}_{#1}}
\newcommand{\pw}[1]{\mathfrak{p}_{#1}}
\newcommand{\GL}{\mathrm{GL}}
\newcommand{\Par}{\operatorname{Par}}
\newcommand{\Iw}{\operatorname{Iw}}
\newcommand{\Tr}{\operatorname{Tr}}
\newcommand{\tr}{\operatorname{tr}}
\newcommand{\lgf}{\operatorname{lg}}
\newcommand{\Forms}{\operatorname{Forms}}
\newcommand{\AllForms}{\operatorname{Forms}_{\kappa}(V)}
\newcommand{\Eint}{\operatorname{End}^{\mathrm{int}}}
\newcommand{\Sys}{\operatorname{Sys}}
\newcommand{\Loc}{\operatorname{Loc}}
\newcommand{\res}{\operatorname{res}}
\newcommand{\Nm}{\operatorname{Nm}}
\newcommand{\Hk}{\mathbb{T}}
\newcommand{\dg}[1]{\langle #1\rangle}
\newcommand{\hp}[2]{\mathsf{h}_{#1,#2}}
\newcommand{\ordp}{\operatorname{ord}^{+}_{\ell}}
\newcommand{\fD}{\mathfrak{D}}
\newcommand{\rr}{\mathfrak{r}}
\newcommand{\bw}{\mathbf{w}}
\newcommand{\bd}{\mathbf{d}}
\newcommand{\one}{\mathbf{1}}
\newcommand{\Fr}{\operatorname{Fr}}
\newcommand{\loc}{\operatorname{loc}}
\newcommand{\Art}{\operatorname{Art}}
\newcommand{\ev}{\operatorname{ev}}
\newcommand{\cA}{\mathcal{A}}
\newcommand{\cS}{\mathcal{S}}
\newcommand{\cJ}{\mathcal{J}}
\newcommand{\cX}{\mathcal{X}}
\newcommand{\cB}{\mathcal{B}}
\newcommand{\cM}{\mathcal{M}}
\newcommand{\tie}{\mathfrak{t}}
\newcommand{\nn}{\mathfrak{n}}
\newcommand{\dif}[1]{\partial_{#1}}
\newcommand{\Yp}{Y_{+}}
\newcommand{\Ym}{Y_{-}}
\newcommand{\rep}[1]{[#1]}
\newcommand{\Hpart}[1]{(#1)_{H}}
\newcommand{\ip}[2]{(#1\mid#2)}
% the full forms of the objects at the level of the auxiliary places
\newcommand{\TQ}{\Hk_{\fD}(\bw_Q)}
\newcommand{\EQ}{\Eint_{\fD}(\bw_Q)}
\newcommand{\FQ}{\Forms_{\fD}(\bw_Q)}
\newcommand{\SQ}{\Sys_{\fD}(\bw_Q)}
\newcommand{\LQ}[1]{\Loc_{\fD}(Q,#1)}
\newcommand{\rQ}{\res_{\fD,Q}}
\newcommand{\NQ}{\Nm_{\fD,Q}}
\newcommand{\inputbox}[1]{\statement{inksoft}{\fontsize{9.8}{13}\selectfont #1}}

\begin{document}
\color{ink}

{\raggedright
{\sffamily\bfseries\fontsize{7.2}{9}\selectfont\color{leancol}\addfontfeatures{LetterSpace=8}READBACK NOTE · CHAPTER VIII · UNIT 137\par}
\vspace{5pt}
{\fontsize{21.5}{25.5}\selectfont\bfseries A Taylor–Wiles prime at a time: ties of tame inertia to Frobenius and the
denominator of the projection\par}
\vspace{5pt}
{\fontsize{12.4}{16}\selectfont\itshape\color{inksoft} One unit of the Lean artifact \ArtifactName, read back into
mathematics\par}}
\vspace{8pt}
{\color{rulegrey}\hrule height 0.6pt}
\vspace{7pt}
\thispagestyle{empty}

\textbf{Summary.} Let $\ell$ be a prime and let $w_j$, \m{j\in J}, be auxiliary (Taylor--Wiles) places of a CM field,
of prime norms \m{q_j\equiv1} modulo~$\ell^\nu$. The finite group \m{\prod_{j}\GL_n(\Fw{w_j})} acts on the algebraic
modular forms of principal level at the~$w_j$ on a definite unitary group of rank~$n$. This note proves two things
about the integral structure of these forms. At one place~$w_j$ there is an integral operator~$F$, supported at
that place, that \emph{ties} the trace of tame inertia minus one to the trace of a polynomial in Frobenius, in the
Galois representations of all Hecke eigensystems at once, modulo~$\ell^\nu$ and up to the factor \m{2\cdot n!}; and
it commutes with the action of \m{\GL_n(\Fw{w_j})} modulo~$\ell^\nu$, up to the factor~$C$ and up to an error that
vanishes after the trace to the base level followed by a Hecke element supported on the eigensystems whose Hecke
polynomial contains a full string of roots. Taking one tie at each place, the top coefficient of Ribet's element, a polynomial in one
variable for each place, gives a Hecke element that carries the projection of an integral form to a chosen set of
eigensystems back into the integral forms, up to a common multiple~$C$ of the orders of the stabilisers
(Theorem~A). The argument is algebra in group rings: a transfer computation in $\GL_n$ of a finite field and a
first-order lemma on weighted sums over a finite group. It rests on Higman's criterion at auxiliary places, on the
chain lemma and on local--global compatibility at pro-$p$ Iwahori level, which are quoted.

\statement{leancol}{\fontsize{9.6}{13.2}\selectfont\textbf{What this note is.} The artifact \ArtifactName{} is a
body of Lean code: the compact edition of a formal proof, due to L.~Alpöge with the AI system Claude, of the Riemann
hypothesis and of its generalisation to Dirichlet $L$-functions. That headline is a claim. The authors report that
every file compiles with Lean v4.35.0 against a pinned version of Mathlib and that Lean's kernel accepted the
proofs; an independent build is recorded in the handout's record of checks when it reports. No mathematician has
read the argument through, and nothing here says that the Riemann hypothesis is proved. This note treats one unit
of the artifact, the folder on one auxiliary prime at a time: all \UnitInNote{} of its \UnitMembers{} declarations.
Theorem~A, Propositions \ref{p:first} and~\ref{p:tie} and Lemmas \ref{l:line} to~\ref{l:charpoly} are proved in
full from the Inputs of Section~\ref{s:setting}: two definitions of another unit of the argument, and results of
the background library, which are stated in boxes and not proved here. Every statement below was compared with its
Lean text, argument by argument. The Lean text read is \PinText; Appendix~A is the concordance of \ConcDecls{}
declarations.}

\sect{The statement}{s:intro}
The twists of Chapter~VIII are chosen so that certain Selmer groups are short. Ribet's method, which belongs to
the unit \emph{Selmer}, turns such a bound into an element of a Hecke algebra with one extra variable~$X_j$ for
each auxiliary place: a relation among traces of Galois representations in which tame inertia at~$w_j$ is
compared with $X_j$ times a polynomial in Frobenius. The present unit is the bridge from that element to modular
forms. In the words of the authors' comment at the head of the file, at one deep auxiliary level the trace to the
old level is deformed in parameters, one for each prime, tying tame inertia to Frobenius, and the top coefficient
of Ribet's Hecke element, a polynomial in them, clears the denominator of the projection to lifts up to a known
factor. The objects are those of Settings \ref{set:forms} to~\ref{set:aux} and of Definitions \ref{d:ord}
to~\ref{d:ribet}.

\statement{arith}{\phantomsection\label{t:A}\textbf{Theorem A} (the top coefficient of Ribet's element clears the
denominator).\enspace\emph{In Settings \ref{set:forms}, \ref{set:twist} and~\ref{set:aux}, let
\m{\beta\colon\SQ\to\Ql} be a function, \m{p\in\Zl[X]} a polynomial with constant coefficient~$p_0$, $\cX$ a subset
of \m{\SQ} and $\pi$ an element of~$\cX$, \m{\varepsilon=(\varepsilon_j)_{j\in J}} a family of polynomials in
\m{\Zl[X]}, \m{L\ge1} an integer, \m{\mu\in\Ql} and \m{A\in[0,\infty]}. Suppose:
\begin{enumerate}[label=\textup{(\arabic*)},leftmargin=2.3em,itemsep=2pt]
\item for every \m{j\in J} and every \m{\xi\in\SQ} that is an eigensystem of base level, has
\m{p(\beta(\xi))\ne0}, and is such that \m{|T(\pi)|<1} for every \m{T\in\TQ} with \m{T(\xi)=0}: the Hecke
polynomial \m{\hp{\xi}{w_j}} has a root string of length~$L$ and ratio~$q_j$ at \m{\omega(\Fr_j)^{-1}} with cofactor
dividing~$\varepsilon_j$ (Definition~\ref{d:string});
\item the element \m{2\cdot n!\cdot(p_0C^2)^{|J|+1}} divides $\ell^\nu$ in~$\Zl$;
\item for every $\Zl$-subalgebra~$\cS$ of the ring of functions from \m{\SQ} to \m{\Ql[X_j:j\in J]} and every
$\cS$-submodule~$\cJ$ of that ring such that \m{(\cS,\cJ)} is a trace ideal for
\m{(\rr,\ \ell^\nu,\ (\tau_j)_{j\in J},\ (\Fr_j)_{j\in J},\ \mathrm{id}_J,\ \varepsilon)}
\textup{(Definition~\ref{d:ideal})}, there is \m{\eta\in\Ql} that is a Ribet coefficient of \m{(\cS,\cJ)}
for~$\beta$ at \m{\pi\in\cX} \textup{(Definition~\ref{d:ribet})} and has \m{\ordp(\mu\eta)\le A}.
\end{enumerate}
Then there is a function \m{s\colon\SQ\to\Ql} with \m{\ordp(\mu\,s(\pi))\le A} and the following property. For
every base-level form \m{\varphi_0\in\Forms_{\fD}(\emptyset)}, every natural number~$\Delta$ such that
\m{\Delta\varphi_0} is integral, and every base-level form \m{\varphi\in\Forms_{\fD}(\emptyset)} such that
\begin{enumerate}[label=\textup{(\roman*)},leftmargin=2.3em,itemsep=2pt]
\item \m{\dg{s'}\varphi=0} for every function \m{s'\colon\SQ\to\Ql} that vanishes on~$\cX$, and
\item \m{\dg{s'}\varphi_0=\dg{s'}\varphi} for every function \m{s'\colon\SQ\to\Ql} that vanishes at all~$\xi$ with
\m{p(\beta(\xi))=0},
\end{enumerate}
the form \m{\Delta\cdot C\cdot\dg{s}\varphi} is integral.}}
\leanline{\lean{mainThm}\sep the formal text has in addition the two hypotheses of Setting~\ref{set:aux}, which
the census does not list (Section~\ref{s:notes})}

Here \m{\dg{s}} multiplies the forms of each eigensystem~$\xi$ by \m{\iota^{-1}(s(\xi))}, and \m{\ordp} is the
$\ell$-adic valuation truncated at~$0$. Condition~(i) says that $\varphi$ is supported on the eigensystems
in~$\cX$, and (ii) that $\varphi$ is the part of~$\varphi_0$ on the eigensystems where \m{p(\beta)} does not
vanish. So $\varphi$ is a projection of~$\varphi_0$, which has denominators in general, and the theorem bounds the
denominator at~$\pi$: after multiplication by~$C$, the element~$s$ clears it, and \m{s(\pi)} is as small as Ribet's
top coefficient~$\eta$ allows. Condition~(1) is the string condition, imposed on the eigensystems~$\xi$ of base
level that are congruent to~$\pi$ in the sense that every Hecke element vanishing at~$\xi$ has absolute value less
than~$1$ at~$\pi$; the unit \emph{Selmer} verifies it where it applies the theorem.

\sect{The setting}{s:setting}
\begin{conventions}\phantomsection\label{c:conventions}
Rings have a unit and need not be commutative; a homomorphism from a group to a ring is a homomorphism into its
multiplicative monoid. For a natural number~$M$ the group \m{\ZZ/M} is written additively; the formal text writes
it multiplicatively. For a finite group and a subgroup~$H$ of it with a homomorphism~$\alpha$ from~$H$ to a
commutative group, \m{\Tr_\alpha} is the transfer of~$\alpha$ to the whole group; its formula is recalled in the
proof of Proposition~\ref{p:first}. In Sections \ref{s:line} and~\ref{s:first} the implicit arguments of the
declarations are introduced in each statement. In the other statements the ambient objects are those of
Setting~\ref{set:forms}: a prime number~$\ell$ with an isomorphism of fields~$\iota$, a CM field~$E$ with a CM type,
embedded in~$\Qbar$, an index set of $n$ elements, a hermitian space~$V$, a weight~$\kappa$, a finite set~$J$, a
natural number~$\nu$, a natural number~$C$ and a base datum~$\fD$; and, where Setting~\ref{set:twist} is cited,
the condition \m{n\ge1}, that $V$ is standard, the Weil data of~$E$, a homomorphism~$\omega$ and a family~$\rr$.
Matrices of size~$n$ have rows and columns numbered from~$1$ to~$n$; the formal text numbers them from~$0$.
\end{conventions}

The three Settings collect the objects of the background library on which the unit is built. They are stated as
the unit's dossier reads them; the notes on the units of the library that own them (their numbers are given)
define them in full.

\begin{setting}[forms, levels, Hecke algebra, auxiliary places; from units 57, 102, 122, 123 and 130 of the
library]\label{set:forms}
\emph{Coefficients.} $\ell$ is a prime number and \m{\iota\colon\CC\to\Ql} a fixed isomorphism of fields, where
$\Ql$ is an algebraic closure of~$\mathbb{Q}_\ell$; $\Zl$ is the ring of integers of~$\Ql$, and \m{|\cdot|} is the
absolute value of~$\Ql$ with \m{|\ell|=\ell^{-1}}. The ring $\Zl$ maps to~$\CC$ by~$\iota^{-1}$, and through this map
it acts on every complex vector space.

\emph{The group.} $E$ is a CM field with a CM type (a complex embedding chosen for each infinite place), embedded
in an algebraic closure~$\Qbar$ of~$\mathbb{Q}$, and \m{\GE=\Gal(\Qbar/E)}. For a finite place~$w$ of~$E$, the
completion is~$\Ew{w}$, its ring of integers~$\Ow{w}$ with maximal ideal~$\pw{w}$, the residue field~$\Fw{w}$ has
$q_w$ elements, and $\bar w$ is the complex conjugate place. $V$ is a positive definite hermitian space over~$E$
with coordinates in an index set of $n$ elements, $\kappa$ is a weight, \m{\Gf} is the group of finite adelic points
of the unitary group of~$V$, which has a standard maximal compact subgroup. \m{\AllForms} is the complex
vector space of the forms of weight~$\kappa$ on~$V$ (algebraic modular forms), on which \m{\Gf} acts by right translation.
At a place \m{w\ne\bar w} the group \m{\GL_n(\Ew{w})} embeds in~\m{\Gf}, and \m{\loc_w} is the resulting
representation of \m{\GL_n(\Ew{w})} on \m{\AllForms}; the \kw{local group algebra} at~$w$ is the $\Zl$-algebra of
operators on \m{\AllForms} generated by the translations by the elements of~\m{\Gf} supported at $w$ and~$\bar w$.

\emph{Base datum.} $C$ is a natural number and $\fD$ a \kw{base datum} for \m{(V,\kappa,C)}: a finite set of finite
places of~$E$; a space of forms that is stable under every operator commuting with the translations by the
components, at these places and their conjugates, of the standard maximal compact subgroup, and on which some open
subgroup acts through a character by its components there; and such that for every \m{\gamma\in\Gf} the order of
the finite group of the rational points that lie in the conjugate by~$\gamma$ of the standard maximal compact
subgroup divides~$C$.

\emph{Levels.} For a family~$\bw$ of finite places of~$E$, \m{\Forms_{\fD}(\bw)} is the space of the forms in the
space of~$\fD$ that are fixed by the principal congruence subgroup at the places of~$\bw$ of the standard maximal
compact subgroup, taken in its components outside the places of~$\fD$ and their conjugates; it is
finite-dimensional. For the empty family this is the space \m{\Forms_{\fD}(\emptyset)} of the \kw{base-level forms},
which is contained in every \m{\Forms_{\fD}(\bw)}. A form is \kw{integral} if $\iota$ maps the coordinates of its
values at the points of~\m{\Gf} that are integral above~$\ell$ into~$\Zl$. \m{\Eint_{\fD}(\bw)} is the $\Zl$-algebra
of the $\CC$-linear endomorphisms of \m{\Forms_{\fD}(\bw)} that map integral forms to integral forms.

\emph{Eigensystems.} An \kw{isotypic space} is a subspace of \m{\AllForms} of the following kind: for some system of
Hecke eigenvalues at all finite places, it consists of the forms that at almost all split places~$w$ are fixed by
\m{\GL_n(\Ow{w})} and are eigenvectors of the spherical Hecke operators at~$w$ with the eigenvalues of the system.
Each isotypic space~$\xi$ has a chosen such system, and \m{\hp{\xi}{w}\in\CC[X]} is the \kw{Hecke polynomial}
of~$\xi$ at~$w$, of degree~$n$, formed from the chosen eigenvalues \m{\theta_{\xi,w,i}\in\CC} at~$w$:
\[\hp{\xi}{w}=\sum_{i=0}^{n}(-1)^{i}\,q_w^{\,i(i-1)/2}\,\theta_{\xi,w,i}\,X^{n-i}.\]
\m{\Sys_{\fD}(\bw)} is the set of the isotypic spaces that contain a non-zero form of \m{\Forms_{\fD}(\bw)}, the
\kw{eigensystems of level~$\bw$}; it is finite, and \m{\Sys_{\fD}(\emptyset)\subseteq\Sys_{\fD}(\bw)} is the set of
the \kw{eigensystems of base level}.

\emph{Diagonal operators and the Hecke algebra.} For a function \m{s\colon\Sys_{\fD}(\bw)\to\Ql}, \m{\dg{s}} is a
chosen endomorphism of \m{\AllForms} that multiplies every form of each \m{\xi\in\Sys_{\fD}(\bw)} by
\m{\iota^{-1}(s(\xi))}; it preserves \m{\Forms_{\fD}(\bw)}. The \kw{Hecke algebra} \m{\Hk_{\fD}(\bw)} is the
$\Zl$-algebra of the functions \m{T\colon\Sys_{\fD}(\bw)\to\Ql} for which the restriction of \m{\dg{T}} to
\m{\Forms_{\fD}(\bw)} lies in \m{\Eint_{\fD}(\bw)}; for such~$T$ the symbol \m{\dg{T}} also denotes this restriction,
and \m{T\mapsto\dg{T}} is a homomorphism of $\Zl$-algebras from \m{\Hk_{\fD}(\bw)} to \m{\Eint_{\fD}(\bw)}.

\emph{Auxiliary places.} $J$ is a finite set, $\nu$ a natural number, and $Q$ is a family of
\kw{Taylor--Wiles places} for \m{(E,J,\ell^\nu)}: for each \m{j\in J} a finite place~$w_j$ of~$E$ whose norm~$q_j$
is a prime number, a place of~$\Qbar$ above~$w_j$, and an element \m{\Fr_j\in\GE} that is a geometric Frobenius
element at that place; the places~$w_j$ are distinct, no conjugate~$\bar w_i$ is among them, and $\ell^\nu$ divides
\m{q_j-1} for every~$j$. \m{\tau_j\in\GE} is the chosen generator of tame inertia at the place above~$w_j$, and
\m{\bw_Q=(w_j)_{j\in J}}. The group \m{\prod_{j\in J}\GL_n(\Fw{w_j})} acts on \m{\FQ} by translation through lifts;
this is a homomorphism
\[\rQ\colon\ \prod_{j\in J}\GL_n(\Fw{w_j})\longrightarrow\EQ,\qquad\text{and}\qquad
\NQ=\sum_{\gamma}\rQ(\gamma)\]
is the sum over the whole group. For \m{j\in J} and \m{\gamma\in\GL_n(\Fw{w_j})}, \m{\gamma^{(j)}} is the element
of the product with $\gamma$ at~$j$ and $1$ elsewhere. \m{\LQ{j}} is the $\Zl$-subalgebra of \m{\EQ} of the
operators that are restrictions of elements of the local group algebra at~$w_j$.
\end{setting}

\begin{setting}[the twisted family of Galois representations; from units 37, 80, 102 and 106]\label{set:twist}
In Setting~\ref{set:forms}, suppose \m{n\ge1}, that $V$ is \kw{standard} (all its signs are~$1$ and its adelic
point is the identity) and that $E$ carries the Weil data of the library. Every isotypic space~$\xi$ has a chosen
homomorphism \m{\varrho_\xi\colon\GE\to\GL_n(\Ql)}: a representation attached to the Hecke polynomials of~$\xi$ if
one exists (Input~3), and arbitrary otherwise. Let \m{\omega\colon\GE\to\Zl} be a homomorphism and let
\m{\rr=(\rr_\xi)_\xi} be a family of homomorphisms \m{\rr_\xi\colon\GE\to\GL_n(\Ql)}, one for each isotypic
space~$\xi$, such that
\[\rr_\xi(\gamma)=\omega(\gamma)\cdot\varrho_\xi(\gamma)\qquad\text{for every isotypic space }\xi\text{ and every }\gamma\in\GE.\]
\end{setting}

\begin{setting}[auxiliary places off level and of trivial character]\label{set:aux}
In Settings \ref{set:forms} and~\ref{set:twist}, suppose:
\begin{enumerate}[label=\textup{(\alph*)},leftmargin=2.3em,itemsep=2pt]
\item for every \m{j\in J} the place $w_j$ is \kw{off level} for~$\fD$ at the base level: neither $w_j$
nor~$\bar w_j$ is a place of~$\fD$, the prime $\ell$ does not lie under~$w_j$, and for every
\m{\xi\in\Sys_{\fD}(\emptyset)} the base-level forms in~$\xi$ are fixed by \m{\GL_n(\Ow{w_j})} and are eigenvectors
of the spherical Hecke operators at~$w_j$ with the chosen eigenvalues \m{\theta_{\xi,w_j,i}} of~$\xi$;
\item \m{\omega(\tau_j)=1} for every \m{j\in J}.
\end{enumerate}
\end{setting}

\begin{definition}[the truncated valuation; unit 92]\label{d:ord}
For \m{z\in\Ql} put \m{\ordp(0)=\infty} and, for \m{z\ne0},
\[\ordp(z)=\max\bigl(0,\,-\log_\ell|z|\bigr)\in[0,\infty).\]
So \m{\ordp(\ell)=1}, and \m{\ordp(z)=0} whenever \m{|z|\ge1}: for instance \m{\ordp(1/\ell)=0}. This is not the
$\ell$-adic valuation, which is negative at~$1/\ell$; the two agree on~$\Zl$. If \m{|z|\le|z'|} then
\m{\ordp(z')\le\ordp(z)}.
\end{definition}

\begin{definition}[root strings; unit 33]\label{d:string}
Let \m{G,G'\in\CC[X]}, let $q$ and~$L$ be natural numbers and \m{\theta\in\CC}. The polynomial~$G$ \kw{has a root
string of length~$L$ and ratio~$q$ at~$\theta$ with cofactor dividing~$G'$} if there is a polynomial
\m{G_1\in\CC[X]} that divides~$G'$ and satisfies
\[G(\theta X)=(X-1)(X-q)\cdots(X-q^{L-1})\cdot G_1.\]
For \m{\theta\ne0} the numbers \m{\theta,\theta q,\dots,\theta q^{L-1}} are then roots of~$G$. Polynomials over~$\Zl$
and elements of~$\Zl$ are used here through the map \m{\iota^{-1}\colon\Zl\to\CC}.
\end{definition}

\begin{definition}[trace ideals; unit 138 of the argument]\label{d:ideal}
Let $I$ be a set, $\cS$ a $\Zl$-subalgebra of the ring of functions from~$I$ to the polynomial ring
\m{\Ql[X_j:j\in J]}, and $\cJ$ an $\cS$-submodule of that ring. Let further a group be given, a family
\m{\rr=(\rr_i)_{i\in I}} of maps from the group to the \m{n\times n} matrices over~$\Ql$, an element \m{y\in\Zl},
two families \m{(\tau_j)} and \m{(\Fr_j)} of elements of the group, indexed by a set with a map
\m{j\mapsto j^{\flat}} to~$J$, and a family \m{\varepsilon=(\varepsilon_j)_{j\in J}} of polynomials in \m{\Zl[X]}.
The pair \m{(\cS,\cJ)} is a \kw{trace ideal for} \m{(\rr,\ y,\ (\tau_j),\ (\Fr_j),\ j\mapsto j^{\flat},\ \varepsilon)} if
\begin{enumerate}[label=\textup{(\alph*)},leftmargin=2.3em,itemsep=2pt]
\item for every \m{j\in J} the constant function with value~$X_j$ lies in~$\cS$;
\item for every element~$\gamma$ of the group and every natural number~$e$, the function that sends \m{i\in I} to
the coefficient of~$X^{e}$ in the characteristic polynomial of \m{\rr_i(\gamma)}, a constant polynomial, lies
in~$\cS$;
\item for every element~$\gamma$ of the group, the function \m{i\mapsto y\cdot\tr\rr_i(\gamma)} lies in~$\cJ$, and
so does, for every index~$j$ of the two families, the function
\[i\ \longmapsto\ \tr\bigl((\rr_i(\tau_j)-1)\,\rr_i(\gamma)\bigr)-X_{j^{\flat}}\cdot
\tr\bigl(\bigl((X-1)\,\varepsilon_{j^{\flat}}\bigr)(\rr_i(\Fr_j))\cdot\rr_i(\gamma)\bigr).\]
\end{enumerate}
In Theorem~A the set~$I$ is \m{\SQ}, the group is~$\GE$, the family~$\rr$ is that of Setting~\ref{set:twist}
restricted to \m{\SQ}, the two families are indexed by~$J$ itself and \m{j^{\flat}=j}. The formal definition, and
that of Definition~\ref{d:ribet}, is made for any commutative ring in the place of~$\Zl$, any commutative algebra
over it in the place of~$\Ql$ and matrices over a ring between the two; only the case stated here is used.
\end{definition}

\begin{definition}[Ribet coefficients; unit 138 of the argument]\label{d:ribet}
Let $I$, $\cS$ and~$\cJ$ be as in Definition~\ref{d:ideal}, let \m{\beta\colon I\to\Ql} be a function, $\cX$ a
subset of~$I$ and \m{\pi\in\cX}. Write \m{\one\in\NN^J} for the exponent of the monomial \m{\prod_{j\in J}X_j}, and
\m{\bd<\one} for the exponents \m{\bd\ne\one} of the monomials that divide it. An element \m{\eta\in\Ql} is a
\kw{Ribet coefficient of \m{(\cS,\cJ)} for~$\beta$ at \m{\pi\in\cX}} if the function~$\beta$, with constant
polynomials as values, lies in~$\cS$ and there is \m{\Phi\in\cS} such that: the product \m{\Phi\cdot\beta} lies
in~$\cJ$; for every \m{i\in\cX}, \m{\beta(i)=0} and the coefficient of the polynomial \m{\Phi(i)} at~$X^{\bd}$
vanishes for every \m{\bd<\one}; and the coefficient of \m{\Phi(\pi)} at \m{\prod_{j\in J}X_j} is~$\eta$.
\end{definition}

\begin{definition}[two subgroups at a place; units 30 and 122]\label{d:par}
Let $w$ be a finite place of a number field~$E$ and \m{n\ge1}. Inside \m{\GL_n(\Ow{w})}, the subgroup
\m{\Par_w} consists of the matrices whose first column lies in~$\pw{w}$ below its first entry, and the
\kw{pro-$p$ Iwahori subgroup} \m{\Iw_w\subseteq\Par_w} of the matrices whose reduction modulo~$\pw{w}$ is upper
triangular with $1$ on the diagonal. More generally a function from \m{\{1,\dots,n\}} to~$\NN$ has a \kw{block
parahoric subgroup}: the matrices in \m{\GL_n(\Ow{w})} whose entry in row~$i$ and column~$i'$ lies in~$\pw{w}$
whenever the function is smaller at~$i'$ than at~$i$; \m{\Par_w} belongs to the function that is $0$ at~$1$ and
$1$ elsewhere, and \m{\GL_n(\Ow{w})} to the function~$0$. For \m{\gamma\in\GL_n(\Ew{w})} and a representation
of \m{\GL_n(\Ew{w})}, \m{U_\gamma} is the Hecke operator of the double coset \m{\Iw_w\gamma\,\Iw_w}: the sum of the
operators of representatives of the cosets \m{\gamma'\,\Iw_w} contained in it. For a uniformiser~$\varpi$
of~$\Ew{w}$, \m{U_\varpi} is this operator for the diagonal matrix with entries \m{\varpi,1,\dots,1}.
\end{definition}

\inputbox{\textbf{Inputs from the background library, 1: the forms.} They are taken for granted, as the dossier
of this unit reads them. $\bw$ is a family of finite places of~$E$ indexed by a finite set.

\vspace{4pt}\textbf{Input 1} (units 102 and 130; elementary).\enspace In Setting~\ref{set:forms}:
\textup{(a)} \m{\Forms_{\fD}(\bw)} is the direct sum of the spaces \m{\xi\cap\Forms_{\fD}(\bw)},
\m{\xi\in\Sys_{\fD}(\bw)}. \textup{(b)} An element of the local group algebra at a place~$w$ maps every isotypic
space into itself, and it commutes with the translation by every element of~\m{\Gf} that is trivial at $w$
and~$\bar w$. \textup{(c)} If $\ell$ does not lie under~$w$, the local group algebra at~$w$ maps integral forms to
integral forms; if \m{w\ne\bar w}, the spherical Hecke operators at~$w$ belong to it. \textup{(d)} The Hecke algebra
\m{\Hk_{\fD}(\bw)} is closed in the space of all functions from the finite set \m{\Sys_{\fD}(\bw)} to~$\Ql$.
\textup{(e)} All but finitely many finite places~$w$ of~$E$ are off level for~$\fD$ at~$\bw$: condition~(a) of
Setting~\ref{set:aux} holds for them, with $\bw$ in the place of the empty family.

\vspace{3pt}\textbf{Input 2} (unit 130; the action at the auxiliary places).\enspace In Setting~\ref{set:forms}:
\textup{(a)} a form \m{\varphi\in\FQ} is of base level if and only if \m{\rQ(\gamma)\varphi=\varphi} for every
\m{\gamma\in\prod_{j}\GL_n(\Fw{w_j})}. \textup{(b)} Every element of \m{\LQ{j}} commutes with \m{\dg{T}} for every
\m{T\in\TQ}. Suppose moreover that the places~$w_j$ are off level (condition~(a) of Setting~\ref{set:aux}) and let
\m{j\in J}. Then: \textup{(c)} for every \m{\gamma\in\GL_n(\Ow{w_j})} with reduction~$\bar\gamma$ and every
\m{\varphi\in\FQ}, \m{\rQ(\bar\gamma^{(j)})\varphi=\loc_{w_j}(\gamma)\varphi}; \textup{(d)} \m{\rQ(\gamma^{(j)})\in\LQ{j}}
for every \m{\gamma\in\GL_n(\Fw{w_j})}; \textup{(e)} if \m{Y\in\LQ{j}} has all its values fixed by \m{\Iw_{w_j}},
then for every \m{\gamma\in\GL_n(\Ew{w_j})} and every \m{G\in\Zl[X]} there is \m{Y'\in\LQ{j}} with
\m{Y'(\varphi)=G(U_\gamma)(Y(\varphi))} for all \m{\varphi\in\FQ}, where $U_\gamma$ is taken in~\m{\loc_{w_j}};
\textup{(f)} \emph{(Higman's criterion up to~$C$)} for every subgroup~$\Gamma$ of \m{\prod_{j}\GL_n(\Fw{w_j})}
there is \m{\Psi\in\EQ} with
\[\sum_{\gamma\in\Gamma}\rQ(\gamma^{-1})\cdot\Psi\cdot\rQ(\gamma)=C\]
(the formal statement is for a finite group with an injective homomorphism to the product, which is the same).
\emph{Reference for \textup{(f)}}, as the catalogue of the
handout gives it: R.~Taylor, A.~Wiles, Ann.\ of Math.\ 141 (1995); L.~Clozel, M.~Harris, R.~Taylor, Publ.\ IHÉS 108
(2008); D.~G.~Higman, Duke Math.\ J.\ 21 (1954); place not confirmed. \emph{Relation} (second reader): classical
freeness assumes trivial stabilisers; the declaration constructs a trace equal to a stabiliser multiple.}

\inputbox{\textbf{Inputs from the background library, continued.}

\vspace{4pt}\textbf{Input 3} (units 3, 5, 80 and 102; attached Galois representations).\enspace In
Setting~\ref{set:twist}, let $\xi$ be an isotypic space that contains a non-zero form. Then \m{\varrho_\xi} is
attached to the Hecke polynomials of~$\xi$ in the sense of the library, and in particular: for every
\m{\gamma\in\GE} and every natural number~$e$, the coefficient of~$X^e$ in \m{\det(X-\varrho_\xi(\gamma))} is the
limit of the $\iota$-images of the coefficients of~$X^e$ in the polynomials \m{\hp{\xi}{w}}, along the filter of
most Frobenius data \m{(w,\cdot,\Fr)} whose Frobenius element tends to~$\gamma$. This filter is proper, its data
live at places of prime norm, and a property that holds for the data at all but finitely many places holds along
it. \emph{Reference}: G.~Chenevier, M.~Harris, Construction of automorphic Galois representations~II, Camb.\
J.\ Math.\ 1 (2013), Theorem~3.2.3. \emph{Relation}: a weaker form (closeness at most Frobenius elements in place
of equality at all unramified split places); the second reader derives it from exact unramified compatibility
after base change and a decomposition of the discrete constituents.

\vspace{3pt}\textbf{Input 4} (unit 102; elementary).\enspace In Setting~\ref{set:forms} the space \m{\AllForms}
carries a positive definite hermitian form \m{\ip{\varphi'}{\varphi}} that is invariant under right translation
by~\m{\Gf}. In particular the operators \m{\loc_w(\gamma)} and, on \m{\FQ}, the operators \m{\rQ(\gamma)} preserve
it.}

\inputbox{\textbf{Inputs from the background library, 2: one split place.} Here $w$ is a finite place of~$E$ with
\m{w\ne\bar w}, $\varpi$ is a uniformiser of~$\Ew{w}$ and \m{u\in\Ew{w}} a unit of~$\Ow{w}$.

\vspace{4pt}\textbf{Input 5} (unit 33; the chain lemma, second part).\enspace In Setting~\ref{set:forms}, let
\m{\varphi'\in\AllForms} be non-zero, fixed by \m{\GL_n(\Ow{w})} and an eigenvector of the spherical Hecke
operators at~$w$ with eigenvalues \m{\theta_i}, and let \m{G=\sum_{i=0}^{n}(-1)^iq_w^{\,i(i-1)/2}\theta_iX^{n-i}}.
Let \m{\theta\in\CC} be non-zero, $e$ a natural number and \m{G'\in\CC[X]}, and put
\m{G''(X)=(\theta X-q_w^{\,e})\cdot G'(\theta X)}. If $G$ has a root string of length \m{e+1} and ratio~$q_w$
at~\m{\theta^{-1}} with cofactor dividing~$G'$, then
\[\ip{\varphi'}{G''(U_\varpi)\,\varphi}=0\qquad\text{for every }\Iw_w\text{-fixed }\varphi\in\AllForms,\]
where $U_\varpi$ is taken in~\m{\loc_w}. The library proves this for every smooth representation of
\m{\GL_n(\Ew{w})} with an invariant positive definite hermitian form. \emph{Reference}, as the catalogue gives it:
M.~Tadić, Ann.\ Sci.\ ENS 19 (1986); J.~Bernstein, Lecture Notes in Math.\ 1041 (1984); place not confirmed, and the
orthogonality stated here has no cited source.

\vspace{3pt}\textbf{Input 6} (unit 30; elementary).\enspace For every function \m{\theta\colon\Ew{w}\to\CC}, the
operator $U_\varpi$ of~\m{\loc_w} maps into itself the space of the \m{\Iw_w}-fixed forms on which every
\m{\gamma\in\Par_w} acts by the scalar \m{\theta(\gamma_{11})}, where \m{\gamma_{11}} is the first entry of~$\gamma$.

\vspace{3pt}\textbf{Input 7} (unit 89; the local Artin map).\enspace Let a place of~$\Qbar$ above a finite
place~$w$ of~$E$ be given, with its decomposition group in~$\GE$, and let \m{\Art_w} be the local Artin map
from~\m{\Ew{w}^{\times}} to the abelianisation of that group, normalised so that uniformisers go to geometric
Frobenius elements. \textup{(a)} If $\tau$ is the chosen generator of tame inertia at that place, there is a unit~$u$
of~$\Ow{w}$ with \m{\Art_w(u)} the class of~$\tau$ whose reduction generates \m{\Fw{w}^{\times}}. \textup{(b)} If
\m{\Fr\in\GE} is a geometric Frobenius element at that place, there is a uniformiser~$\varpi$ with \m{\Art_w(\varpi)}
the class of~$\Fr$. \emph{Reference}: J.-P.~Serre, Local Fields (1979), Chapters XII to~XIV. \emph{Relation}:
weaker forms (existence statements in the image of local reciprocity).}

\inputbox{\textbf{Inputs from the background library, 2, continued.} The place~$w$, the uniformiser~$\varpi$ and
the unit~$u$ are as before.

\vspace{4pt}\textbf{Input 8} (unit 32; local--global compatibility at tame level).\enspace In
Setting~\ref{set:twist}, suppose that $\ell$ does not lie under~$w$, and let a place of~$\Qbar$ above~$w$ be given
together with elements $\Fr$ and~$\tau$ of its decomposition group such that \m{\Art_w(\varpi)} is the class
of~$\Fr$, \m{\Art_w(u)} is the class of~$\tau$, $\Fr$ is a geometric Frobenius element and $\tau$ lies in the
inertia group. Let $\xi$ be an isotypic space and \m{\varphi\in\xi}.

\textup{(a)} Suppose \m{\varphi\ne0} is \m{\Iw_w}-fixed and every~$\gamma$ in the block parahoric subgroup of some
function acts on~$\varphi$ by a scalar \m{\theta(\gamma)\in\CC}. Let \m{\theta_i} be the value of~$\theta$ at the
diagonal matrix with $u$ at the place~$i$ and $1$ elsewhere, and suppose that \m{\theta_i=\theta_{i'}} only if the
function has the same value at $i$ and~$i'$. Then every polynomial over~$\Ql$ that vanishes at all
\m{\iota(\theta_i)} vanishes at the matrix \m{\varrho_\xi(\tau)}; the matrices \m{\varrho_\xi(\Fr)} and
\m{\varrho_\xi(\tau)} commute; and there are characters \m{\theta'_1,\dots,\theta'_n} of~\m{\Ew{w}^{\times}}, forming an
exponent of~\m{\loc_w} with an exponent functional that does not vanish on~$\varphi$, such that for every natural
number~$e$ the polynomial \m{\prod_{i}(X-\iota(\theta'_i(\varpi)\theta_i^{\,e}))} vanishes at
\m{\varrho_\xi(\Fr)\varrho_\xi(\tau)^{e}}.

\textup{(b)} Suppose $\varphi$ is \m{\Iw_w}-fixed and every \m{\gamma\in\Par_w} acts on~$\varphi$ by the scalar
\m{\theta(\gamma_{11})}, for a function \m{\theta\colon\Ew{w}\to\CC} with \m{\theta(u)\ne0} and
\m{\theta(u)^{n!}\ne1}. Then for every \m{G\in\CC[X]} and every \m{n\times n} matrix~$y$ over~$\Ql$,
\[\iota^{-1}\tr\bigl((\varrho_\xi(\tau)-1)\,y\bigr)\cdot G(U_\varpi)\varphi
=\iota^{-1}\tr\bigl(G^{\iota}(\varrho_\xi(\Fr))\,(\varrho_\xi(\tau)-1)\,y\bigr)\cdot\varphi,\]
where \m{G^{\iota}\in\Ql[X]} is the image of~$G$ under~$\iota$. \emph{Reference}, as the catalogue gives it:
R.~Taylor, T.~Yoshida, J.\ Amer.\ Math.\ Soc.\ 20 (2007); M.~Harris, R.~Taylor (2001); place not confirmed.
\emph{Relation} (second reader): for~(a), the published comparison away from~$\ell$ needs a local type and exponent
calculation to recover these simultaneous matrix identities; for~(b), a matrix-weighted identity for a specific
Hecke operator requires additional parahoric linear algebra. Both comparisons are open.}

Inputs 1(a) and~4 are used in the following form. Two endomorphisms of \m{\Forms_{\fD}(\bw)} that agree on
\m{\xi\cap\Forms_{\fD}(\bw)} for every \m{\xi\in\Sys_{\fD}(\bw)} are equal; and since \m{\dg{T}} multiplies the
forms of~$\xi$ by \m{\iota^{-1}(T(\xi))}, a form \m{\varphi\in\Forms_{\fD}(\bw)} with components
\m{\varphi_\xi\in\xi} has \m{\dg{T}\varphi=\sum_\xi\iota^{-1}(T(\xi))\varphi_\xi}.

\sect{The line stabiliser in $\GL_n$ of a finite field}{s:line}
\begin{lemma}[a character of the line stabiliser whose transfer is a function of the determinant]\label{l:line}
Let $E$ be a number field, $w$ a finite place of~$E$, \m{n\ge1} an integer and $M$ a natural number that divides
\m{q_w-1}. Let \m{u\in\Ew{w}} be a unit of~$\Ow{w}$ whose reduction~$\bar u$ generates the group
\m{\Fw{w}^{\times}}. Then there are a subgroup~$H$ of \m{\GL_n(\Fw{w})}, a homomorphism \m{\alpha\colon H\to\ZZ/M},
two functions \m{\lgf,\delta\colon\Ew{w}\to\ZZ/M} and an element \m{\gamma_0\in\GL_n(\Fw{w})} with the following
properties, in which \m{\bar\gamma\in\GL_n(\Fw{w})} is the reduction of \m{\gamma\in\GL_n(\Ow{w})}, with first
entry~\m{\gamma_{11}}.
\begin{enumerate}[label=\textup{(\alph*)},leftmargin=2.3em,itemsep=2pt]
\item For every \m{\gamma\in\GL_n(\Ow{w})}: \m{\bar\gamma\in H} if and only if \m{\gamma\in\Par_w}; if
\m{\bar\gamma\in H} then \m{\alpha(\bar\gamma)=\lgf(\gamma_{11})}; \m{\Tr_\alpha(\bar\gamma)=\delta(\det\gamma)}; and
if \m{\gamma\in\Iw_w} then \m{\Tr_\alpha(\bar\gamma)=0} and \m{\alpha(\bar\gamma)=0}.
\item \m{\lgf(u)=1}, \m{\Tr_\alpha(\gamma_0)=\delta(u)}, and \m{n!\cdot\Tr_\alpha(\gamma_0)=n!} in \m{\ZZ/M}.
\end{enumerate}
\end{lemma}
\begin{proof}
Write \m{\Gamma=\GL_n(\Fw{w})}. Let $H$ be the stabiliser in~$\Gamma$ of the first coordinate line: the matrices~$h$
whose first column vanishes below its first entry~\m{h_{11}}. For \m{h,h'\in H} the first entry of \m{hh'} is
\m{h_{11}h'_{11}}, so \m{h\mapsto h_{11}} is a homomorphism from~$H$ to \m{\Fw{w}^{\times}}. This group is cyclic of
order \m{q_w-1} and generated by~$\bar u$; since $M$ divides \m{q_w-1} there is a homomorphism
\m{\lgf_0\colon\Fw{w}^{\times}\to\ZZ/M} with \m{\lgf_0(\bar u)=1}. Put \m{\alpha(h)=\lgf_0(h_{11})}, and for
\m{z\in\Fw{w}^{\times}} let \m{\delta_0(z)=\Tr_\alpha(\diag(z,1,\dots,1))}; this is a homomorphism from
\m{\Fw{w}^{\times}} to \m{\ZZ/M}.

\emph{The transfer is \m{\delta_0\circ\det}.} Both \m{\Tr_\alpha} and \m{\delta_0\circ\det} are homomorphisms
from~$\Gamma$ to a commutative group, and $\Gamma$ is generated by the diagonal matrices and the transvections:
the matrices that differ from the identity matrix by one entry \m{z\in\Fw{w}}, in a row~$i$ and a column \m{i'\ne i}.
A transvection has determinant~$1$. For fixed \m{i\ne i'} the transvection with entry \m{z+z'} is the product of
those with entries $z$ and~$z'$, so the transfer of the transvection with entry~$z$ is a homomorphism in~$z$
from the additive group of~$\Fw{w}$, of order~$q_w$, to \m{\ZZ/M}; so $q_w$ times each of its values
is~$0$, and since \m{q_w\equiv1} modulo~$M$ the values themselves are~$0$. A diagonal matrix
\m{\diag(z_1,\dots,z_n)} is the product of the matrices with $z_i$ at the place~$i$ and $1$ elsewhere. Such a matrix
is conjugate, by a permutation matrix, to \m{\diag(z_i,1,\dots,1)}, and the transfer, a homomorphism to a
commutative group, is constant on conjugacy classes. So \m{\Tr_\alpha(\diag(z_1,\dots,z_n))=\sum_i\delta_0(z_i)
=\delta_0(z_1\cdots z_n)}. Hence \m{\Tr_\alpha=\delta_0\circ\det} on~$\Gamma$.

\emph{The functions and the element.} For a unit~$z$ of~$\Ow{w}$ with reduction~$\bar z$ put
\m{\lgf(z)=\lgf_0(\bar z)} and \m{\delta(z)=\delta_0(\bar z)}; at the other elements of~$\Ew{w}$ the two functions
take arbitrary values, which play no role. Let \m{\gamma_0=\diag(\bar u,1,\dots,1)}.

\emph{Property (a).} Let \m{\gamma\in\GL_n(\Ow{w})}. Its reduction lies in~$H$ exactly when the first column
of~$\gamma$ lies in~$\pw{w}$ below the first entry, that is, when \m{\gamma\in\Par_w}. In that case \m{\gamma_{11}} is
a unit and \m{\alpha(\bar\gamma)=\lgf_0(\bar\gamma_{11})=\lgf(\gamma_{11})}. The determinant commutes with
reduction, so \m{\Tr_\alpha(\bar\gamma)=\delta_0(\det\bar\gamma)=\delta(\det\gamma)}. If \m{\gamma\in\Iw_w}, then
$\bar\gamma$ is upper triangular with $1$ on the diagonal: it lies in~$H$, its determinant is~$1$ and its first
entry is~$1$, whence \m{\Tr_\alpha(\bar\gamma)=\delta_0(1)=0} and \m{\alpha(\bar\gamma)=\lgf_0(1)=0}.

\emph{Property (b).} \m{\lgf(u)=\lgf_0(\bar u)=1} and \m{\Tr_\alpha(\gamma_0)=\delta_0(\bar u)=\delta(u)} by the
definitions. For the last equation consider the scalar matrix \m{\bar u\cdot1_n}. It is central in~$\Gamma$ and
lies in~$H$, so for every representative~$\varsigma$ of a coset \m{\varsigma H} one has
\m{(\bar u\cdot1_n)\varsigma=\varsigma\,(\bar u\cdot1_n)}, and formula~(\ref{e:transfer}) for the transfer, in
the proof of Proposition~\ref{p:first}, gives \m{\Tr_\alpha(\bar u\cdot1_n)=[\Gamma:H]\cdot\alpha(\bar u\cdot1_n)=[\Gamma:H]}. The
index of~$H$ is the number of lines in~\m{\Fw{w}^{\,n}}, which is \m{1+q_w+\dots+q_w^{\,n-1}\equiv n} modulo~$M$. On the
other hand \m{\Tr_\alpha(\bar u\cdot1_n)=\delta_0(\bar u^{\,n})=n\cdot\delta_0(\bar u)}. So
\m{n\cdot\Tr_\alpha(\gamma_0)=n} in \m{\ZZ/M}, and multiplication by \m{(n-1)!} gives
\m{n!\cdot\Tr_\alpha(\gamma_0)=n!}.
\end{proof}
\leanline{\lean{lineStab}\sep the formal text takes the unit as an element of~\m{\Ew{w}^{\times}} with a proof
that it lies in the unit group of~$\Ow{w}$}

\sect{Weighted sums over a finite group and the first-order lemma}{s:first}
\begin{definition}[weighted sums]\label{d:wsum}
Let $\Gamma$ be a finite group, $R$ a ring and $M$ a natural number. For homomorphisms \m{\rho\colon\Gamma\to R} and
\m{\chi\colon\Gamma\to\ZZ/M} and a function \m{f\colon\ZZ/M\to\ZZ}, the \kw{weighted sum} is
\[\weightedSum{\rho}{\chi}{f}=\sum_{\gamma\in\Gamma}f(\chi(\gamma))\cdot\rho(\gamma)\ \in R.\]
For fixed $\rho$ and~$\chi$ it is additive in~$f$, and \m{\weightedSum{\rho}{\chi}{\,\cdot\,}} is this additive map
from the group of the functions \m{\ZZ/M\to\ZZ} to~$R$. The constant function with value~$1$ is written~$1$, and
for \m{z'\in\ZZ/M} or \m{z'\in\NN} the function \m{z\mapsto f(z-z')} is written \m{f(\,\cdot-z')}.
\end{definition}
\leanline{\lean{weightedSum}}

\begin{lemma}[translation of the weights]\label{l:shift}
Let $\Gamma$ be a finite group, $R$ a ring, $M$ a natural number, \m{\rho\colon\Gamma\to R} and
\m{\chi\colon\Gamma\to\ZZ/M} homomorphisms, \m{f\colon\ZZ/M\to\ZZ} a function and \m{\gamma\in\Gamma}. Then
\[\rho(\gamma)\cdot\weightedSum{\rho}{\chi}{f}=\weightedSum{\rho}{\chi}{f(\,\cdot-\chi(\gamma))}.\]
\end{lemma}
\begin{proof}
\m{\rho(\gamma)\sum_{\gamma'}f(\chi(\gamma'))\rho(\gamma')=\sum_{\gamma'}f(\chi(\gamma'))\rho(\gamma\gamma')}, and the
substitution \m{\gamma''=\gamma\gamma'} turns this into
\m{\sum_{\gamma''}f(\chi(\gamma'')-\chi(\gamma))\rho(\gamma'')}.
\end{proof}
\leanline{\lean{mulWeightedSum}}

\begin{lemma}[weights that agree modulo~$M$]\label{l:modeq}
Let $R$ be a ring, $S$ a ring, $M$ a natural number and \m{\sigma\colon R\to S} a ring homomorphism with
\m{M\cdot1_S=0}. Let $I$ be a set, $B$ an additive map from the group of the functions \m{I\to\ZZ} to~$R$, and
\m{f,f'\colon I\to\ZZ} two functions with \m{f'(i)\equiv f(i)} modulo~$M$ for every \m{i\in I}. Then
\m{\sigma(B(f))=\sigma(B(f'))}.
\end{lemma}
\begin{proof}
There is a function \m{f''\colon I\to\ZZ} with \m{f-f'=M\cdot f''}, so
\[\sigma(B(f))-\sigma(B(f'))=M\cdot\sigma(B(f''))=0.\]
\end{proof}
\leanline{\lean{mapEqWeights}}

In the next statement \m{\dif{e}f=f(\,\cdot-e)-f} is the difference of a function \m{f\colon\ZZ/M\to\ZZ} with
step~$e$, and \m{\rho|_H} is the restriction of~$\rho$ to a subgroup~$H$.

\begin{proposition}[the first-order lemma]\label{p:first}
Let $\Gamma$ be a finite group, $R$ a ring, $M$ a natural number, \m{\rho\colon\Gamma\to R} a homomorphism, $H$ a
subgroup of~$\Gamma$, $S$ a ring and \m{\sigma\colon R\to S} a ring homomorphism with \m{M\cdot1_S=0}. Let
\m{\alpha\colon H\to\ZZ/M} be a homomorphism and $B$ an additive map from the group of the functions
\m{\ZZ/M\to\ZZ} to~$R$ such that
\[B(f)\cdot\rho(h)=B\bigl(f(\,\cdot-\alpha(h))\bigr)\qquad\text{for all }f\text{ and all }h\in H.\]
Let $\cA$ be a subring of~$R$ that contains \m{\rho(\gamma)} for every \m{\gamma\in\Gamma} and \m{B(f)} for
every~$f$, and let \m{(c_i)_{i\in I}} be a family of elements of~$R$, indexed by a set~$I$, each of which commutes
with every element of~$\cA$. Let \m{\Psi,\Pi\in R}, \m{\gamma_0\in\Gamma} and \m{C,e\in\NN} be such that
\begin{enumerate}[label=\textup{(H\arabic*)},leftmargin=3em,itemsep=2pt]
\item \m{\sum_{h\in H}\rho(h^{-1})\,\Psi\,\rho(h)=C};
\item \m{\rho(h)\cdot B(1)=B(1)} for every \m{h\in H};
\item \m{\Pi\cdot B(1)=0};
\item \m{\Pi\cdot\rho(\gamma)=\Pi} for every \m{\gamma\in\Gamma};
\item multiplication by the order \m{|\Gamma|} is injective on~$R$;
\item \m{e\cdot\Tr_\alpha(\gamma_0)=e} in \m{\ZZ/M}.
\end{enumerate}
Then there is \m{F\in\cA} with the following two properties.
\begin{enumerate}[label=\textup{(\roman*)},leftmargin=2.3em,itemsep=2pt]
\item For every \m{\gamma\in\Gamma} there is \m{Z\in R} with
\[\sigma\bigl(C\cdot F\cdot(\rho(\gamma)-1)\bigr)=\sigma\Bigl(\Bigl(\sum_{\gamma'\in\Gamma}\rho(\gamma')\Bigr)\cdot
Z\Bigr)\qquad\text{and}\qquad\Pi\cdot Z=0.\]
\item For all \m{i,i'\in I} and \m{y\in R}: if
\begin{gather*}
B(\dif{e}f)\cdot c_i=\weightedSum{\rho|_H}{\alpha}{\dif{e}f}\cdot c_{i'}\quad\text{for every }f,\qquad\text{and}\\
\weightedSum{\rho}{\Tr_\alpha}{f}\cdot c_{i'}=(\rho(\gamma_0)-1)\cdot\weightedSum{\rho}{\Tr_\alpha}{f}\cdot y
\quad\text{for every }f,
\end{gather*}
then
\[2e\cdot\sigma(F)\,\sigma(c_i)=2e\cdot\sigma\Bigl(\sum_{\gamma\in\Gamma}\rho(\gamma)\Bigr)\,\sigma(y).\]
\end{enumerate}
\end{proposition}
\begin{proof}
\emph{Coset sums.} Fix a set \m{\Sigma\subseteq\Gamma} of representatives of the left cosets \m{\gamma H}. For
\m{\gamma\in\Gamma} let \m{\rep{\gamma}\in\Sigma} be the representative of \m{\gamma H} and
\m{\Hpart{\gamma}=\rep{\gamma}^{-1}\gamma\in H}, so that \m{\gamma=\rep{\gamma}\cdot\Hpart{\gamma}}. Put
\[\Yp=\sum_{\varsigma\in\Sigma}\rho(\varsigma),\qquad\Ym=\sum_{\varsigma\in\Sigma}\rho(\varsigma^{-1}),\qquad
\Nm=\sum_{\gamma\in\Gamma}\rho(\gamma),\]
and let \m{v\colon\ZZ/M\to\ZZ} send a residue to its representative of least absolute value, so that
\m{v(z)\equiv z} modulo~$M$. The transfer of~$\alpha$ is, by definition and independently of~$\Sigma$,
\begin{equation}\label{e:transfer}
\Tr_\alpha(\gamma)=\sum_{\varsigma\in\Sigma}\alpha\bigl(\Hpart{\gamma\varsigma}\bigr)\qquad(\gamma\in\Gamma).
\end{equation}
For a function \m{f_0\colon H\to\ZZ} and \m{\gamma'\in\Gamma},
\begin{equation}\label{e:key}
\Yp\cdot\Bigl(\sum_{h\in H}f_0(h)\rho(h)\Bigr)\cdot\rho(\gamma'^{-1})
=\sum_{\gamma\in\Gamma}f_0\bigl(\Hpart{\gamma\gamma'}\bigr)\,\rho(\gamma).
\end{equation}
Indeed the left side is the sum of \m{f_0(h)\rho(\varsigma h\gamma'^{-1})} over \m{(\varsigma,h)\in\Sigma\times H}, and
\m{(\varsigma,h)\mapsto\varsigma h\gamma'^{-1}} is a bijection from \m{\Sigma\times H} to~$\Gamma$ whose inverse sends
$\gamma$ to \m{(\rep{\gamma\gamma'},\Hpart{\gamma\gamma'})}. For \m{f_0=1} and \m{\gamma'=1} this reads
\m{\Yp\sum_{h\in H}\rho(h)=\Nm}. Hence, by (H2) and then (H1) with $h^{-1}$ in the place of~$h$,
\begin{equation}\label{e:higman}
\Nm\cdot\Psi\cdot B(1)=\Yp\sum_{h\in H}\rho(h)\,\Psi\,\rho(h^{-1})\,B(1)=C\cdot\Yp\cdot B(1).
\end{equation}

\emph{The operator.} Let \m{F=\Yp\cdot B(-v)\cdot\Ym}. It lies in~$\cA$, because $\cA$ is a subring that contains
the \m{\rho(\gamma)} and \m{B(-v)}.

\emph{Property (i).} Let \m{\gamma\in\Gamma}. The map \m{\varsigma\mapsto\rep{\gamma\varsigma}} is a permutation
of~$\Sigma$, and \m{\rep{\gamma\varsigma}^{-1}\gamma=\Hpart{\gamma\varsigma}\,\varsigma^{-1}}; therefore
\m{\Ym\,\rho(\gamma)=\sum_{\varsigma}\rho(\Hpart{\gamma\varsigma})\,\rho(\varsigma^{-1})}, and by the hypothesis on~$B$
\[F\cdot\rho(\gamma)=\Yp\sum_{\varsigma\in\Sigma}B\bigl(-v(\,\cdot-\alpha(\Hpart{\gamma\varsigma}))\bigr)\,\rho(\varsigma^{-1}).\]
For \m{z'\in\ZZ/M} the functions \m{-v(\,\cdot-z')} and \m{-v+v(z')\cdot1} agree modulo~$M$ at every point, so
\m{\sigma(B(-v(\,\cdot-z')))=\sigma(B(-v))+v(z')\,\sigma(B(1))} by Lemma~\ref{l:modeq}. Hence
\[\sigma\bigl(F\cdot(\rho(\gamma)-1)\bigr)=\sigma\Bigl(\Yp\cdot B(1)\cdot\sum_{\varsigma\in\Sigma}
v\bigl(\alpha(\Hpart{\gamma\varsigma})\bigr)\,\rho(\varsigma^{-1})\Bigr).\]
Multiply by~$C$ and use~(\ref{e:higman}): property~(i) holds with
\[Z=\Psi\cdot B(1)\cdot\sum_{\varsigma\in\Sigma}v\bigl(\alpha(\Hpart{\gamma\varsigma})\bigr)\,\rho(\varsigma^{-1}).\]
It remains to see that \m{\Pi\cdot\Psi\cdot B(1)=0}. By~(H4), \m{\Pi\cdot\Nm=|\Gamma|\cdot\Pi} and
\m{\Pi\cdot\Yp=[\Gamma:H]\cdot\Pi}. Multiplying (\ref{e:higman}) by~$\Pi$ on the left gives
\m{|\Gamma|\cdot\Pi\Psi B(1)=C\,[\Gamma:H]\cdot\Pi B(1)=0} by~(H3), and (H5) gives \m{\Pi\Psi B(1)=0}.

\emph{Property (ii).} Let $i$, $i'$ and~$y$ satisfy the two hypotheses of~(ii). First, by (\ref{e:key}) for
\m{f_0=v\circ\alpha} and \m{\gamma'=\varsigma}, summed over \m{\varsigma\in\Sigma},
\[\Yp\cdot\weightedSum{\rho|_H}{\alpha}{v}\cdot\Ym=\sum_{\gamma\in\Gamma}\Bigl(\sum_{\varsigma\in\Sigma}
v\bigl(\alpha(\Hpart{\gamma\varsigma})\bigr)\Bigr)\rho(\gamma),\]
and by (\ref{e:transfer}) the inner sum is congruent to \m{v(\Tr_\alpha(\gamma))} modulo~$M$. So
Lemma~\ref{l:modeq}, applied to functions on~$\Gamma$, gives
\begin{equation}\label{e:conj}
\sigma\bigl(\Yp\cdot\weightedSum{\rho|_H}{\alpha}{v}\cdot\Ym\bigr)=\sigma\bigl(\weightedSum{\rho}{\Tr_\alpha}{v}\bigr):
\end{equation}
conjugating by the representatives sums the weights over the cosets. Secondly, take in the first hypothesis the
function \m{f=-v\cdot(v+e)}. Since \m{v(z-e)\equiv v(z)-e} modulo~$M$,
\[\dif{e}f(z)\equiv-(v(z)-e)\,v(z)+v(z)\,(v(z)+e)=2e\cdot v(z)\pmod M\]
for every~$z$: the function~$f$ is a discrete antiderivative of \m{2e\cdot v}, and this is where the factor~$2$
of the statement comes from. So, by Lemma~\ref{l:modeq},
\[\sigma(B(\dif{e}f))=2e\cdot\sigma(B(v))\qquad\text{and}\qquad
\sigma\bigl(\weightedSum{\rho|_H}{\alpha}{\dif{e}f}\bigr)=2e\cdot\sigma\bigl(\weightedSum{\rho|_H}{\alpha}{v}\bigr).\]
Multiply the first hypothesis by~$\Yp$ on the left and by~$\Ym$ on the right, apply~$\sigma$,
and move $c_i$ and~$c_{i'}$ to the right of \m{\Ym\in\cA}, with which they commute. By the definition of~$F$ and
by~(\ref{e:conj}) this gives
\[-2e\cdot\sigma(F)\,\sigma(c_i)=2e\cdot\sigma\bigl(\weightedSum{\rho}{\Tr_\alpha}{v}\bigr)\,\sigma(c_{i'})
=\sigma\bigl(\weightedSum{\rho}{\Tr_\alpha}{2e\cdot v}\cdot c_{i'}\bigr).\]
Thirdly, by the second hypothesis for the function \m{2e\cdot v} and by Lemma~\ref{l:shift},
\[\weightedSum{\rho}{\Tr_\alpha}{2e\cdot v}\cdot c_{i'}
=\weightedSum{\rho}{\Tr_\alpha}{2e\cdot\bigl(v(\,\cdot-\Tr_\alpha(\gamma_0))-v\bigr)}\cdot y,\]
and \m{2e\,(v(z-\Tr_\alpha(\gamma_0))-v(z))\equiv-2e\cdot\Tr_\alpha(\gamma_0)\equiv-2e} modulo~$M$ by~(H6). By
Lemma~\ref{l:modeq} once more, $\sigma$ of the right side is
\m{\sigma(\weightedSum{\rho}{\Tr_\alpha}{-2e\cdot1})\,\sigma(y)=-2e\cdot\sigma(\Nm)\,\sigma(y)}. The two displayed
equations together give \m{2e\cdot\sigma(F)\sigma(c_i)=2e\cdot\sigma(\Nm)\sigma(y)}.
\end{proof}
\leanline{\lean{firstOrder}\sep the formal text writes the family as a map from an index type and takes
$\Psi$, $\Pi$, $\gamma_0$, $C$ and~$e$ as implicit arguments}

\sect{The tie at one auxiliary place}{s:tie}
\begin{lemma}[testing on eigenvectors]\label{l:test}
In Setting~\ref{set:forms}, let \m{j\in J}, let \m{T,T'\in\TQ}, let $\Gamma$ be a finite group,
\m{\rho\colon\Gamma\to\EQ} a homomorphism with \m{\rho(\gamma)\in\LQ{j}} for every \m{\gamma\in\Gamma},
\m{\chi\colon\Gamma\to\ZZ/\ell^\nu} a homomorphism, \m{f\colon\ZZ/\ell^\nu\to\ZZ} a function and \m{Y\in\EQ}. Let
\m{e\in\ZZ/\ell^\nu} and \m{\theta\in\CC}, let $\cB$ be a $\CC$-linear map from \m{\FQ} to \m{\AllForms} such that
\m{Y(\varphi)=\cB\bigl(\weightedSum{\rho}{\chi}{f}\varphi\bigr)} for every \m{\varphi\in\FQ}, and let
\m{\tilde f\colon\ZZ/\ell^\nu\to\CC} be a function with \m{\tilde f(z-e)-\theta\cdot\tilde f(z)=f(z)} for
every~$z$. Suppose that
\[\iota^{-1}(T(\xi))\cdot\varphi=\iota^{-1}(T'(\xi))\cdot\cB(\varphi)\]
for every \m{\xi\in\SQ}, every additive character \m{\psi\colon\ZZ/\ell^\nu\to\CC^{\times}} and every non-zero
\m{\varphi\in\xi\cap\FQ} such that \m{\rho(\gamma)\varphi=\psi(\chi(\gamma))\cdot\varphi} for all
\m{\gamma\in\Gamma} and \m{\psi(e)\ne\theta}. Then
\[\weightedSum{\rho}{\chi}{f}\cdot\dg{T}=Y\cdot\dg{T'}\qquad\text{in }\EQ.\]
\end{lemma}
\begin{proof}
By Input~1(a) it suffices to compare the two sides on a form \m{\varphi\in\xi\cap\FQ}, for each \m{\xi\in\SQ}. For
a function \m{f'\colon\ZZ/\ell^\nu\to\CC} put
\[\varphi[f']=\sum_{\gamma\in\Gamma}\bigl(f'(\chi(\gamma)-e)-\theta\cdot f'(\chi(\gamma))\bigr)\cdot\rho(\gamma)\varphi.\]
This is linear in~$f'$, and \m{\varphi[\tilde f]=\weightedSum{\rho}{\chi}{f}\varphi}. The left side of the
assertion sends $\varphi$ to \m{\iota^{-1}(T(\xi))\cdot\varphi[\tilde f]} and the right side to
\m{\iota^{-1}(T'(\xi))\cdot Y(\varphi)=\iota^{-1}(T'(\xi))\cdot\cB(\varphi[\tilde f])}. So it suffices to prove
\[\iota^{-1}(T(\xi))\cdot\varphi[f']=\iota^{-1}(T'(\xi))\cdot\cB(\varphi[f'])\]
for all~$f'$, and by linearity for the additive characters \m{f'=\psi}, which form a basis of the functions. Now
\[\varphi[\psi]=\bigl(\psi(-e)-\theta\bigr)\cdot\sum_{\gamma\in\Gamma}\psi(\chi(\gamma))\,\rho(\gamma)\varphi,\qquad
\rho(\gamma')\,\varphi[\psi]=\psi^{-1}(\chi(\gamma'))\cdot\varphi[\psi]\quad(\gamma'\in\Gamma),\]
the second by the substitution \m{\gamma\mapsto\gamma'\gamma} in the sum, and \m{\varphi[\psi]\in\xi} by
Input~1(b), since each \m{\rho(\gamma)} is the restriction of an element of the local group algebra at~$w_j$. If
\m{\varphi[\psi]=0}, both sides of the equation for \m{f'=\psi} vanish. This is the case whenever
\m{\psi^{-1}(e)=\theta}; it may also be the case otherwise. If \m{\varphi[\psi]\ne0}, then
\m{\psi^{-1}(e)=\psi(-e)\ne\theta}, and the hypothesis, applied to $\xi$, to the character~$\psi^{-1}$ and to the
form \m{\varphi[\psi]}, is the equation.
\end{proof}
\leanline{\lean{weightedSumExt}\sep the formal text does not use that $V$ is standard, the Weil data or
\m{n\ge1}}

\begin{lemma}[coefficients of characteristic polynomials lie in the Hecke algebra]\label{l:charpoly}
In Settings \ref{set:forms} and~\ref{set:twist}, let $\bw$ be a family of finite places of~$E$ indexed by~$J$, let
\m{\gamma\in\GE} and let $e$ be a natural number. Then the function on \m{\Sys_{\fD}(\bw)} that sends~$\xi$ to the
coefficient of~$X^e$ in \m{\det(X-\rr_\xi(\gamma))} belongs to \m{\Hk_{\fD}(\bw)}.
\end{lemma}
\begin{proof}
Since \m{\rr_\xi(\gamma)=\omega(\gamma)\varrho_\xi(\gamma)}, the coefficient is \m{\omega(\gamma)^{n-e}} times the
coefficient of~$X^e$ in \m{\det(X-\varrho_\xi(\gamma))} for \m{e\le n}, and it is $0$ for \m{e>n}. The Hecke algebra
is a $\Zl$-algebra and \m{\omega(\gamma)\in\Zl}, so it suffices to treat~\m{\varrho}. Every
\m{\xi\in\Sys_{\fD}(\bw)} contains a non-zero form, so by Input~3 the function in question is the limit, in the
space of functions on the finite set \m{\Sys_{\fD}(\bw)}, of the functions
\m{\xi\mapsto\iota(\text{coefficient of }X^e\text{ in }\hp{\xi}{w})} along a proper filter on Frobenius data at
places~$w$ of prime norm. The Hecke algebra is closed (Input~1(d)), and all but finitely many places are off level
for~$\fD$ at~$\bw$ (Input~1(e)). So it suffices to show that the function belongs to \m{\Hk_{\fD}(\bw)} for a
place~$w$ of prime norm that is off level. Such a place has \m{w\ne\bar w}. The coefficient of~$X^e$ in
\m{\hp{\xi}{w}} is an integer, depending only on $w$, $n$ and~$e$, times the chosen eigenvalue
\m{\theta_{\xi,w,n-e}} of the spherical Hecke operator of index \m{n-e} at~$w$. Because $w$ is off level, the forms
in \m{\xi\cap\Forms_{\fD}(\bw)} are eigenvectors of that operator with this eigenvalue, for every~$\xi$; so by
Input~1(a) the diagonal operator of the function \m{\xi\mapsto\iota(\theta_{\xi,w,n-e})} agrees on
\m{\Forms_{\fD}(\bw)} with the spherical Hecke operator. That operator lies in the local group algebra at~$w$ and
$\ell$ does not lie under~$w$, so it maps integral forms to integral forms (Input~1(c)). Hence the function
belongs to \m{\Hk_{\fD}(\bw)}, and so does its multiple by an integer.
\end{proof}
\leanline{\lean{charpolyMem}}

\begin{definition}[trace elements]\label{d:trace}
In Settings \ref{set:forms} and~\ref{set:twist}, let $\bw$ be a family of finite places of~$E$ indexed by~$J$. For a
polynomial \m{G\in\Zl[X]} and elements \m{\gamma',\gamma\in\GE}, the \kw{trace element}
\m{\traceElem{\rr}{G}{\gamma'}{\gamma}} is the function
\[\traceElem{\rr}{G}{\gamma'}{\gamma}\colon\ \Sys_{\fD}(\bw)\to\Ql,\qquad
\xi\longmapsto\tr\bigl(G(\rr_\xi(\gamma'))\cdot\rr_\xi(\gamma)\bigr).\]
It belongs to \m{\Hk_{\fD}(\bw)}: the matrix \m{G(\rr_\xi(\gamma'))\rr_\xi(\gamma)} is a combination, with the
coefficients of~$G$, of the matrices \m{\rr_\xi(\gamma'^{\,i}\gamma)}, and the trace of \m{\rr_\xi(\gamma'')} is
minus the coefficient of~\m{X^{n-1}} in \m{\det(X-\rr_\xi(\gamma''))}, to which Lemma~\ref{l:charpoly} applies. The
index~$\rr$ stands for the family~$\rr$ together with the relation
\m{\rr_\xi(\gamma)=\omega(\gamma)\cdot\varrho_\xi(\gamma)} of Setting~\ref{set:twist}, which is the first argument of
the formal definition. For instance \m{\traceElem{\rr}{X-1}{\tau_j}{\gamma}} sends~$\xi$ to
\m{\tr((\rr_\xi(\tau_j)-1)\rr_\xi(\gamma))}, and \m{\traceElem{\rr}{1}{1}{\gamma}} sends~$\xi$ to
\m{\tr\rr_\xi(\gamma)}.
\end{definition}
\leanline{\lean{traceElem}}

\begin{proposition}[the tie of tame inertia to Frobenius at one auxiliary place]\label{p:tie}
In Settings \ref{set:forms}, \ref{set:twist} and~\ref{set:aux}, let \m{j\in J}, let \m{\varepsilon\in\Zl[X]} be a
polynomial, \m{L\ge1} an integer and \m{T\in\TQ}. Suppose that for every \m{\xi\in\SQ} that is an eigensystem of
base level and has \m{T(\xi)\ne0}, the Hecke polynomial \m{\hp{\xi}{w_j}} has a root string of length~$L$ and
ratio~$q_j$ at \m{\omega(\Fr_j)^{-1}} with cofactor dividing~$\varepsilon$ \textup{(Definition~\ref{d:string})}. Let $S$
be a ring and \m{\sigma\colon\EQ\to S} a ring homomorphism with \m{\ell^\nu\cdot1_S=0}. Put
\[\rho_j(\gamma)=\rQ(\gamma^{(j)})\quad\text{for }\gamma\in\GL_n(\Fw{w_j}),\qquad\text{and}\qquad
\Nm_j=\sum_{\gamma\in\GL_n(\Fw{w_j})}\rho_j(\gamma).\]
Then there is \m{F\in\LQ{j}} with the following two properties.
\begin{enumerate}[label=\textup{(\roman*)},leftmargin=2.3em,itemsep=2pt]
\item For every \m{\gamma\in\GL_n(\Fw{w_j})} there is \m{Z\in\EQ} with
\[\sigma\bigl(C\cdot F\cdot(\rho_j(\gamma)-1)\bigr)=\sigma(\Nm_j\cdot Z)\qquad\text{and}\qquad
\dg{T}\cdot\NQ\cdot Z=0\ \text{ in }\EQ.\]
\item For every \m{\gamma\in\GE}, the element \m{y=\sigma\bigl(\dg{\traceElem{\rr}{X-1}{\tau_j}{\gamma}}\bigr)}
of~$S$ satisfies
\[\sigma(\Nm_j)\cdot y=0\qquad\text{and}\qquad
2\cdot n!\cdot\sigma(F)\cdot y=2\cdot n!\cdot\sigma(\Nm_j)\cdot
\sigma\bigl(\dg{\traceElem{\rr}{(X-1)\,\varepsilon}{\Fr_j}{\gamma}}\bigr).\]
\end{enumerate}
The trace elements are taken at the level~$\bw_Q$.
\end{proposition}
\leanline{\lean{existsTie}\sep the formal statement ends at line~103 of its file: the census cuts it at the
binding of \m{\rho_j}}

In words: modulo~$\ell^\nu$ and up to the factor \m{2\cdot n!}, the operator~$F$ at~$w_j$ times the trace of
\m{(\tau_j-1)\gamma} is the sum~\m{\Nm_j} over \m{\GL_n(\Fw{w_j})}, the trace to the level without~$w_j$, times the
trace of \m{(\Fr_j-1)\,\varepsilon(\Fr_j)\,\gamma}. The tame generator and the Frobenius element enter two different
trace elements: the polynomial \m{X-1} is evaluated at~$\tau_j$, and \m{(X-1)\varepsilon} at~$\Fr_j$. Property~(i)
measures how far $F$ is from commuting with the action of \m{\GL_n(\Fw{w_j})}: up to the factor~$C$ and
modulo~$\ell^\nu$, the operator \m{F\cdot(\rho_j(\gamma)-1)} is a multiple of~\m{\Nm_j} by an operator that
\m{\dg{T}\cdot\NQ} annihilates. This is the form in which the authors' comment, ``ties hold on forms and weak lifts
do not move'', enters the proof of Theorem~A (Step~8).

\begin{proof}
Write \m{w=w_j} and \m{q=q_j}. By Input~7 there are a unit~$u$ of~$\Ow{w}$ whose reduction generates
\m{\Fw{w}^{\times}} and whose image under~\m{\Art_w} is the class of~$\tau_j$, and a uniformiser~$\varpi$ of~$\Ew{w}$
whose image is the class of~$\Fr_j$. Since $\ell^\nu$ divides \m{q-1}, Lemma~\ref{l:line} applies to
\m{(E,w,n,\ell^\nu,u)}; it gives a subgroup~$H$ of \m{\GL_n(\Fw{w})}, a homomorphism \m{\alpha\colon H\to\ZZ/\ell^\nu},
functions $\lgf$ and~$\delta$ and an element~$\gamma_0$. Let
\[\varepsilon^{+}=(X-q^{L-1})\cdot\varepsilon\qquad\text{and}\qquad
\varepsilon^{\sharp}(X)=\varepsilon^{+}\bigl(\omega(\Fr_j)\cdot X\bigr),\]
polynomials over~$\Zl$. By Setting~\ref{set:aux}(a) the place~$w$ is off level, so $\ell$ does not lie under~$w$ and
Input~2(c) to~(f) are available; and \m{w\ne\bar w}.

\emph{Three remarks on fixed vectors.}
\begin{itemize}[leftmargin=2.6em,itemsep=2pt]
\item[\textup{(R1)}] For \m{\gamma\in\GL_n(\Ow{w})} and every form $\varphi$ in \m{\FQ} one has
\m{\loc_w(\gamma)\varphi=\rho_j(\bar\gamma)\varphi} (Input~2(c)).
\item[\textup{(R2)}] A form \m{\varphi\in\FQ} that is fixed by \m{\rho_j(h)} for all \m{h\in H} with
\m{\alpha(h)=0} and \m{\Tr_\alpha(h)=0} is fixed by \m{\Iw_w}: an element of \m{\Iw_w} lies in \m{\Par_w}, its
reduction~$h$ lies in~$H$ and has \m{\alpha(h)=0=\Tr_\alpha(h)} by Lemma~\ref{l:line}(a), and (R1) applies.
\item[\textup{(R3)}] For every \m{f\colon\ZZ/\ell^\nu\to\ZZ} and \m{\varphi\in\FQ} the form
\m{\weightedSum{\rho_j|_H}{\alpha}{f}\varphi} is fixed by \m{\Iw_w}: for \m{h\in H} with \m{\alpha(h)=0},
Lemma~\ref{l:shift} gives
\[\rho_j(h)\cdot\weightedSum{\rho_j|_H}{\alpha}{f}=\weightedSum{\rho_j|_H}{\alpha}{f},\]
and (R2) applies.
\end{itemize}

\emph{The operators \m{B(f)}.} The operator \m{\weightedSum{\rho_j|_H}{\alpha}{f}} lies in \m{\LQ{j}} by
Input~2(d), and its values are fixed by \m{\Iw_w} by~(R3). So Input~2(e) gives, for every~$f$, an operator
\[B(f)\in\LQ{j}\qquad\text{with}\qquad
B(f)\varphi=\varepsilon^{\sharp}(U_\varpi)\bigl(\weightedSum{\rho_j|_H}{\alpha}{f}\varphi\bigr)\quad(\varphi\in\FQ),\]
and $B$ is additive in~$f$. We apply Proposition~\ref{p:first} to the group \m{\Gamma=\GL_n(\Fw{w})}, the ring
\m{R=\EQ}, the modulus \m{M=\ell^\nu}, the homomorphism~\m{\rho_j}, the subgroup~$H$, the ring~$S$ with~$\sigma$,
the homomorphism~$\alpha$, the map~$B$, the subring \m{\cA=\LQ{j}}, the family \m{(\dg{T'})_{T'\in\TQ}} indexed by
the Hecke algebra, a Higman operator~$\Psi$ for the subgroup \m{\{h^{(j)}:h\in H\}} (Input~2(f)), the element
\m{\Pi=\dg{T}\cdot\NQ}, the element~$\gamma_0$, the integer~$C$ and \m{e=n!}. Its hypotheses hold:
\begin{itemize}[leftmargin=1.4em,itemsep=2pt]
\item \emph{The hypothesis on~$B$.} For \m{h\in H}, the substitution \m{h'\mapsto h'h} in the sum gives
\[\weightedSum{\rho_j|_H}{\alpha}{f}\cdot\rho_j(h)=\weightedSum{\rho_j|_H}{\alpha}{f(\,\cdot-\alpha(h))};\]
apply \m{\varepsilon^{\sharp}(U_\varpi)}.
\item \emph{The subring and the family.} \m{\rho_j(\gamma)\in\LQ{j}} by Input~2(d), \m{B(f)\in\LQ{j}} by
construction, and every \m{\dg{T'}} commutes with \m{\LQ{j}} by Input~2(b).
\item \emph{\textup{(H1)}} is the property of~$\Psi$, and \emph{\textup{(H6)}} is Lemma~\ref{l:line}(b).
\item \emph{\textup{(H2)}.} Let \m{h\in H}, \m{\varphi\in\FQ} and
\m{\varphi'=\weightedSum{\rho_j|_H}{\alpha}{1}\varphi}. By (R3) and Lemma~\ref{l:shift} the form~$\varphi'$ is fixed
by \m{\Iw_w} and by every \m{\rho_j(h')}, \m{h'\in H}; by (R1) and Lemma~\ref{l:line}(a) it is fixed by~\m{\Par_w}.
By Input~6 for the constant function~$1$, the operator \m{U_\varpi} preserves the space of such forms, hence so does
\m{\varepsilon^{\sharp}(U_\varpi)}. So \m{B(1)\varphi} is fixed by~\m{\Par_w}; as reduction maps \m{\Par_w} onto~$H$,
(R1) gives \m{\rho_j(h)B(1)\varphi=B(1)\varphi}.
\item \emph{\textup{(H4)}.} \m{\NQ\cdot\rQ(\gamma^{(j)})=\NQ}, a sum over a group.
\item \emph{\textup{(H5)}.} $R$ is a ring of endomorphisms of a complex vector space.
\end{itemize}
Hypothesis~(H3), \m{\dg{T}\cdot\NQ\cdot B(1)=0}, is the heart of the matter and is proved next. Granting it, the
proposition gives \m{F\in\LQ{j}} whose property~(i) is property~(i) above.

\emph{Proof of \textup{(H3)}: the chain lemma.} Let \m{\varphi\in\FQ} and \m{\varphi_1=\dg{T}\,\NQ\,B(1)\varphi}; we
show \m{\ip{\varphi_1}{\varphi_1}=0} for the invariant form of Input~4. The operator \m{\dg{T}} commutes with
\m{B(1)} and with every \m{\rQ(\gamma)}, a product of operators in the algebras \m{\LQ{j'}} (Input~2(b)
and~(d)). So \m{\varphi_1=\NQ\,B(1)\dg{T}\varphi} and also \m{\varphi_1=\dg{T}\varphi_2} with
\m{\varphi_2=\NQ\,B(1)\varphi}. The form~$\varphi_2$ is fixed by every \m{\rQ(\gamma)}, so it is of base level
(Input~2(a)), and by Input~1(a) for the empty family it is a sum of forms
\m{\varphi'\in\xi\cap\Forms_{\fD}(\emptyset)}, \m{\xi\in\Sys_{\fD}(\emptyset)}. For such a form
\m{\dg{T}\varphi'=\iota^{-1}(T(\xi))\varphi'}. So it suffices to show \m{\ip{\varphi'}{\varphi_1}=0} for every
\m{\xi\in\Sys_{\fD}(\emptyset)} with \m{T(\xi)\ne0} and every non-zero \m{\varphi'\in\xi\cap\Forms_{\fD}(\emptyset)}.
Such a~$\varphi'$ is fixed by every \m{\rQ(\gamma)}, and these operators preserve the form, hence
\begin{gather*}
\ip{\varphi'}{\varphi_1}=\sum_{\gamma}\ip{\varphi'}{\rQ(\gamma)B(1)\dg{T}\varphi}
=\Bigl|\prod_{j'\in J}\GL_n(\Fw{w_{j'}})\Bigr|\cdot\ip{\varphi'}{\varepsilon^{\sharp}(U_\varpi)\varphi_3},\\
\text{where}\qquad\varphi_3=\weightedSum{\rho_j|_H}{\alpha}{1}\dg{T}\varphi,
\end{gather*}
and $\varphi_3$ is fixed by \m{\Iw_w} by~(R3). By Setting~\ref{set:aux}(a) the form~$\varphi'$ is fixed by
\m{\GL_n(\Ow{w})} and is an eigenvector of the spherical Hecke operators at~$w$ with the chosen eigenvalues
of~$\xi$, so its Hecke polynomial is \m{\hp{\xi}{w}}. By the hypothesis of the proposition, \m{\hp{\xi}{w}} has a
root string of length \m{L=(L-1)+1} and ratio~$q$ at \m{\omega(\Fr_j)^{-1}} with cofactor dividing~$\varepsilon$, and
\m{\omega(\Fr_j)} is a unit of~$\Zl$, hence not zero. The chain lemma (Input~5, with \m{\theta=\omega(\Fr_j)},
\m{e=L-1} and \m{G'=\varepsilon}) gives \m{\ip{\varphi'}{G''(U_\varpi)\varphi_3}=0} for
\m{G''(X)=(\omega(\Fr_j)X-q^{L-1})\cdot\varepsilon(\omega(\Fr_j)X)=\varepsilon^{\sharp}(X)}. This proves~(H3).

\emph{Property (ii).} Fix \m{\gamma\in\GE}. For \m{G\in\Zl[X]} put
\[T_G=\traceElem{\rr}{G}{\Fr_j}{\gamma}\qquad\text{and}\qquad
T^{\tau}_G=\traceElem{\rr}{G}{\Fr_j}{\tau_j\gamma}-\traceElem{\rr}{G}{\Fr_j}{\gamma},\]
elements of \m{\TQ}; so \m{T^{\tau}_G(\xi)=\tr\bigl(G(\rr_\xi(\Fr_j))\,(\rr_\xi(\tau_j)-1)\,\rr_\xi(\gamma)\bigr)}, and
\m{T^{\tau}_1=\traceElem{\rr}{X-1}{\tau_j}{\gamma}}. Recall \m{\rr_\xi(\tau_j)=\varrho_\xi(\tau_j)} by
Setting~\ref{set:aux}(b).

\emph{Claim A: where \m{\GL_n} acts through the determinant, the tame generator acts as a diamond operator.} For
every \m{G\in\Zl[X]} and every \m{f\colon\ZZ/\ell^\nu\to\ZZ},
\begin{equation}\label{e:diamond}
\weightedSum{\rho_j}{\Tr_\alpha}{f}\cdot\dg{T^{\tau}_G}=(\rho_j(\gamma_0)-1)\cdot\weightedSum{\rho_j}{\Tr_\alpha}{f}\cdot\dg{T_G}.
\end{equation}
This is Lemma~\ref{l:test} for the group \m{\GL_n(\Fw{w})} with \m{\rho_j} and \m{\chi=\Tr_\alpha}, with \m{e=0},
\m{\theta=0} and \m{\tilde f=f}, with \m{\cB=\rho_j(\gamma_0)-1} and
\m{Y=(\rho_j(\gamma_0)-1)\weightedSum{\rho_j}{\Tr_\alpha}{f}}, and with the Hecke elements \m{T^{\tau}_G} and~\m{T_G}.
Its hypothesis is checked as follows. Let \m{\xi\in\SQ}, let $\psi$ be an additive character and
\m{\varphi\in\xi\cap\FQ} a non-zero form with \m{\rho_j(\gamma')\varphi=\psi(\Tr_\alpha(\gamma'))\varphi} for all
\m{\gamma'\in\GL_n(\Fw{w})}; the condition \m{\psi(0)\ne0} is empty. By~(R2) the form~$\varphi$ is fixed by
\m{\Iw_w}, and by (R1) and Lemma~\ref{l:line}(a) every \m{\gamma'\in\GL_n(\Ow{w})} acts on it by the scalar
\m{\psi(\delta(\det\gamma'))}, which is \m{\psi(\delta(u))} at each diagonal matrix with one entry~$u$ and the
others~$1$. Input~8(a), for the function~$0$, whose block parahoric subgroup is \m{\GL_n(\Ow{w})}, shows that the
polynomial \m{X-\iota(\psi(\delta(u)))} vanishes at \m{\varrho_\xi(\tau_j)}: this matrix is the scalar
\m{\iota(\psi(\delta(u)))}. Hence \m{T^{\tau}_G(\xi)=\bigl(\iota(\psi(\delta(u)))-1\bigr)\cdot T_G(\xi)}. On the other
hand \m{(\rho_j(\gamma_0)-1)\varphi=(\psi(\Tr_\alpha(\gamma_0))-1)\varphi=(\psi(\delta(u))-1)\varphi} by
Lemma~\ref{l:line}(b). So \m{\iota^{-1}(T^{\tau}_G(\xi))\varphi=\iota^{-1}(T_G(\xi))\cdot(\rho_j(\gamma_0)-1)\varphi},
as required.

\emph{The first assertion of \textup{(ii)}.} Take \m{G=1} and \m{f=1} in~(\ref{e:diamond}). Then
\m{\weightedSum{\rho_j}{\Tr_\alpha}{1}=\Nm_j} and \m{(\rho_j(\gamma_0)-1)\Nm_j=0}, so \m{\Nm_j\cdot\dg{T^{\tau}_1}=0}
already in \m{\EQ}, and \m{\sigma(\Nm_j)\cdot y=0}.

\emph{Claim B: on the eigenvectors of the line stabiliser, the trace against \m{\tau_j-1} turns
\m{U_\varpi} into Frobenius.} For every \m{f\colon\ZZ/\ell^\nu\to\ZZ},
\begin{equation}\label{e:frob}
B(\dif{n!}f)\cdot\dg{T^{\tau}_1}=\weightedSum{\rho_j|_H}{\alpha}{\dif{n!}f}\cdot\dg{T^{\tau}_{\varepsilon^{+}}}.
\end{equation}
This is Lemma~\ref{l:test} for the group~$H$ with \m{\rho_j|_H} and \m{\chi=\alpha}, with \m{e=n!}, \m{\theta=1},
the integral function \m{\dif{n!}f} and \m{\tilde f=f}, with $\cB$ the restriction of
\m{\varepsilon^{\sharp}(U_\varpi)} to \m{\FQ} and \m{Y=B(\dif{n!}f)}, and with the Hecke elements
\m{T^{\tau}_{\varepsilon^{+}}} and~\m{T^{\tau}_1}. To check its hypothesis let \m{\xi\in\SQ}, let $\psi$ be an additive
character with \m{\psi(n!)\ne1} and \m{\varphi\in\xi\cap\FQ} a non-zero form with
\m{\rho_j(h)\varphi=\psi(\alpha(h))\varphi} for all \m{h\in H}. By (R2) the form~$\varphi$ is fixed by \m{\Iw_w}, and
by (R1) and Lemma~\ref{l:line}(a) every \m{\gamma'\in\Par_w} acts on it by the scalar
\m{\psi(\lgf(\gamma'_{11}))}. The function \m{\theta=\psi\circ\lgf} on~$\Ew{w}$ has \m{\theta(u)=\psi(1)\ne0} and
\m{\theta(u)^{n!}=\psi(n!)\ne1}. Input~8(b), for the image of~$\varepsilon^{\sharp}$ in \m{\CC[X]} and the matrix
\m{\rr_\xi(\gamma)}, gives
\[\iota^{-1}\tr\bigl((\varrho_\xi(\tau_j)-1)\rr_\xi(\gamma)\bigr)\cdot\varepsilon^{\sharp}(U_\varpi)\varphi
=\iota^{-1}\tr\bigl(\varepsilon^{\sharp}(\varrho_\xi(\Fr_j))\,(\varrho_\xi(\tau_j)-1)\,\rr_\xi(\gamma)\bigr)\cdot\varphi.\]
Here
\[\varepsilon^{\sharp}(\varrho_\xi(\Fr_j))=\varepsilon^{+}\bigl(\omega(\Fr_j)\varrho_\xi(\Fr_j)\bigr)=\varepsilon^{+}(\rr_\xi(\Fr_j)),\]
so the equation reads
\[\iota^{-1}(T^{\tau}_1(\xi))\cdot\cB(\varphi)=\iota^{-1}(T^{\tau}_{\varepsilon^{+}}(\xi))\cdot\varphi,\]
as required.

\emph{The tie.} Apply property~(ii) of Proposition~\ref{p:first} to the indices \m{T^{\tau}_1} and
\m{T^{\tau}_{\varepsilon^{+}}} of the family and to the element \m{\dg{T_{\varepsilon^{+}}}} of~$R$: its first hypothesis
is~(\ref{e:frob}) and its second is (\ref{e:diamond}) for \m{G=\varepsilon^{+}}. It gives
\[2\cdot n!\cdot\sigma(F)\cdot\sigma(\dg{T^{\tau}_1})=2\cdot n!\cdot\sigma(\Nm_j)\cdot\sigma(\dg{T_{\varepsilon^{+}}}).\]
Finally \m{T_{\varepsilon^{+}}=\traceElem{\rr}{X\varepsilon}{\Fr_j}{\gamma}-q^{L-1}\cdot\traceElem{\rr}{\varepsilon}{\Fr_j}{\gamma}},
and \m{q^{L-1}\equiv1} modulo~$\ell^\nu$ because $\ell^\nu$ divides \m{q-1}; as $\ell^\nu$ vanishes in~$S$,
\m{\sigma(\dg{T_{\varepsilon^{+}}})=\sigma(\dg{\traceElem{\rr}{(X-1)\varepsilon}{\Fr_j}{\gamma}})}.
\end{proof}

\sect{Proof of Theorem A}{s:proofA}
\begin{proof}[Proof of Theorem A]
All trace elements are taken at the level~$\bw_Q$. A multi-index is an element $\bd$ of~\m{\NN^J}, with
components~\m{\bd_j}. Multi-indices are ordered componentwise, \m{|\bd|} is the sum of the components, \m{\one}
has all components~$1$, \m{\one_j} has the component~$1$ at~$j$ and $0$ elsewhere, and
\m{X^{\bd}=\prod_jX_j^{\bd_j}}.

\emph{Step 1: separating elements.} Call \m{\xi\in\SQ} \kw{separated from~$\pi$} if there is \m{T\in\TQ} with
\m{T(\xi)=0} and \m{|T(\pi)|\ge1}; choose such an element~$T_\xi$ for each separated~$\xi$, and let \m{T_{*}} be
the product of these finitely many elements. Then \m{|T_{*}(\pi)|\ge1}, and if \m{T_{*}(\xi)\ne0} then $\xi$ is not
separated from~$\pi$: every \m{T\in\TQ} with \m{T(\xi)=0} has \m{|T(\pi)|<1}.

\emph{Step 2: Ribet's element.} Let \m{\ev} be the map from the polynomial ring \m{\TQ[X_j:j\in J]} to the ring of
functions from \m{\SQ} to \m{\Ql[X_j:j\in J]} that evaluates the coefficients at each eigensystem; it is an
injective homomorphism of $\Zl$-algebras. Let $\cS$ be its image, let $\tie$ be the ideal of \m{\TQ[X_j:j\in J]}
generated by the elements
\[\ell^\nu\cdot\traceElem{\rr}{1}{1}{\gamma}\qquad\text{and}\qquad
\traceElem{\rr}{X-1}{\tau_j}{\gamma}-X_j\cdot\traceElem{\rr}{(X-1)\,\varepsilon_j}{\Fr_j}{\gamma}
\qquad(\gamma\in\GE,\ j\in J),\]
and let \m{\cJ=\ev(\tie)}, an $\cS$-submodule. The pair \m{(\cS,\cJ)} is a trace ideal for
\[(\rr,\ \ell^\nu,\ (\tau_j)_{j\in J},\ (\Fr_j)_{j\in J},\ \mathrm{id}_J,\ \varepsilon):\]
the variables lie in~$\cS$, the coefficients of the characteristic polynomials lie in~$\cS$ by Lemma~\ref{l:charpoly} at the
level~$\bw_Q$, and the functions of Definition~\ref{d:ideal}(c) are the images under~\m{\ev} of the generators
of~$\tie$. By hypothesis~(3) there is a Ribet coefficient \m{\eta\in\Ql} of \m{(\cS,\cJ)} for~$\beta$ at
\m{\pi\in\cX} with \m{\ordp(\mu\eta)\le A}. Unfolding Definition~\ref{d:ribet}: the function~$\beta$ is the image
under~\m{\ev} of a polynomial, hence, comparing constant terms, \m{\beta\in\TQ}; and there is a polynomial
\m{\Phi=\sum_{\bd}\Phi_{\bd}X^{\bd}} with coefficients \m{\Phi_{\bd}\in\TQ} such that, \m{\ev} being injective,
\begin{equation}\label{e:ribet}
\beta\cdot\Phi\in\tie;\qquad\beta(\xi)=0\ \text{ and }\ \Phi_{\bd}(\xi)=0\quad(\xi\in\cX,\ \bd<\one);\qquad
\Phi_{\one}(\pi)=\eta.
\end{equation}

\emph{Step 3: the function~$s$.} Put
\[T_0=T_{*}\cdot p(\beta)\in\TQ\qquad\text{and}\qquad s=T_{*}^{\,|J|+1}\cdot\Phi_{\one}.\]
Then \m{|\mu\,s(\pi)|\ge|\mu\,\eta|} by Step~1, so \m{\ordp(\mu\,s(\pi))\le\ordp(\mu\eta)\le A}
(Definition~\ref{d:ord}). It remains to prove the property of~$s$. Fix $\varphi_0$, $\Delta$ and~$\varphi$ as in
the statement; both forms lie in \m{\FQ}.

\emph{Step 4: constants.} By hypothesis~(2) there is \m{\upsilon\in\Zl} with
\begin{equation}\label{e:nu}
\ell^\nu=2\cdot n!\cdot(p_0C^2)^{|J|+1}\cdot\upsilon.
\end{equation}
Since \m{p(\beta)-p_0} is a multiple of~$\beta$ in \m{\TQ}, there is \m{T'\in\TQ} with
\begin{equation}\label{e:power}
(C^2\,T_0)^{|J|+1}=(p_0C^2)^{|J|+1}\cdot T_{*}^{\,|J|+1}+\beta\cdot T'.
\end{equation}
By Input~2(f) for the whole group \m{\prod_{j}\GL_n(\Fw{w_j})} there is \m{\Psi\in\EQ} with
\[\sum_{\gamma}\rQ(\gamma^{-1})\,\Psi\,\rQ(\gamma)=C.\]

\emph{Step 5: the module that $\varphi_0$ spans over the integral operators, modulo~$\ell^\nu$.} Let
\[\nn=\bigl\{\Upsilon\in\EQ:\ \Upsilon(\varphi_0)=\ell^\nu\cdot\Upsilon'(\varphi_0)\text{ for some }\Upsilon'\in\EQ\bigr\}.\]
This is a left ideal of \m{\EQ} that contains \m{\ell^\nu\cdot\EQ}. Let \m{\cM=\EQ/\nn}, a left module over \m{\EQ}
on which $\ell^\nu$ vanishes, let \m{\sigma\colon\EQ\to S} be the action, with $S$ the ring of the $\Zl$-linear
endomorphisms of~$\cM$, and write \m{[\Upsilon]\in\cM} for the class of~$\Upsilon$. (The map
\m{\Upsilon\mapsto\Upsilon(\varphi_0)} identifies $\cM$ with the lattice \m{\EQ\cdot\varphi_0} modulo~$\ell^\nu$.) We
claim: for every \m{T''\in\TQ} that is a multiple of~$T_0$ in~\m{\TQ},
\begin{equation}\label{e:crit}
\sigma(\dg{T''})\,\sigma(\NQ)\,[\Psi]=0\quad\Longleftrightarrow\quad\text{there is }\Upsilon\in\EQ\text{ with }
\ell^\nu\cdot\Upsilon(\varphi_0)=C\cdot\dg{T''}\varphi.
\end{equation}
Indeed the left side says that \m{\dg{T''}\,\NQ\,\Psi\in\nn}, that is,
\m{\dg{T''}\NQ\Psi(\varphi_0)=\ell^\nu\Upsilon(\varphi_0)} for some~$\Upsilon$. The form~$\varphi_0$ is of base
level, hence fixed by every \m{\rQ(\gamma)} (Input~2(a)), so
\m{\NQ\Psi(\varphi_0)=\sum_\gamma\rQ(\gamma)\Psi\rQ(\gamma^{-1})\varphi_0=C\cdot\varphi_0} by the property
of~$\Psi$. And \m{\dg{T''}\varphi_0=\dg{T''}\varphi} by hypothesis~(ii), because $T''$ is a multiple of
\m{p(\beta)} and so vanishes wherever \m{p(\beta)} does.

\emph{Step 6: the ties.} Let \m{j\in J}. Proposition~\ref{p:tie} applies to~$j$, to the polynomial~$\varepsilon_j$,
to~$L$, to the Hecke element~$T_0$ and to \m{\sigma\colon\EQ\to S}: if \m{\xi\in\SQ} is of base level and
\m{T_0(\xi)\ne0}, then \m{p(\beta(\xi))\ne0} and \m{T_{*}(\xi)\ne0}, so by Step~1 every $T$ with \m{T(\xi)=0} has
\m{|T(\pi)|<1}, and hypothesis~(1) gives the root string. We obtain \m{F_j\in\LQ{j}} such that, with
\m{\rho_j} and~\m{\Nm_j} as in the proposition:
\begin{enumerate}[label=\textup{(\roman*)},leftmargin=2.3em,itemsep=2pt]
\item for every \m{\gamma\in\GL_n(\Fw{w_j})} there is \m{Z\in\EQ} with
\[\sigma\bigl(C\cdot F_j\cdot(\rho_j(\gamma)-1)\bigr)=\sigma(\Nm_j\cdot Z)\qquad\text{and}\qquad
\dg{T_0}\cdot\NQ\cdot Z=0;\]
\item for every \m{\gamma\in\GE}, the two Hecke elements
\[T_1=\traceElem{\rr}{X-1}{\tau_j}{\gamma}\qquad\text{and}\qquad
T_2=\traceElem{\rr}{(X-1)\,\varepsilon_j}{\Fr_j}{\gamma}\]
satisfy
\[\sigma(\Nm_j)\,\sigma(\dg{T_1})=0\qquad\text{and}\qquad
2\,n!\cdot\sigma(F_j)\,\sigma(\dg{T_1})=2\,n!\cdot\sigma(\Nm_j)\,\sigma(\dg{T_2}).\]
\end{enumerate}

\emph{Step 7: the operators~\m{\Xi_{\bd}}.} The places~$w_j$ are distinct and none is conjugate to another, so the
images under~$\sigma$ of \m{\LQ{j}} and \m{\LQ{j'}} commute for \m{j\ne j'} (Input~1(b)). For \m{\bd\in\NN^J} let
\[\Xi_{\bd}=\prod_{j\in J}\sigma\bigl(S_j(\bd)\bigr)\in S,\qquad\text{where}\quad
S_j(\bd)=\begin{cases}F_j&\text{if }\bd_j=0,\\ \Nm_j&\text{if }\bd_j=1,\\ 0&\text{if }\bd_j\ge2.\end{cases}\]
In the authors' words, \m{\Xi_{\bd}} is the coefficient complementary to~$\bd$ of the family
\m{\prod_j(\Nm_j+X_jF_j)}. Four properties:
\begin{enumerate}[label=\textup{(\alph*)},leftmargin=2.3em,itemsep=2pt]
\item \m{\Xi_{\bd}} commutes with \m{\sigma(\dg{T})} for every \m{T\in\TQ} (Input~2(b));
\item \m{\Xi_{\one}=\sigma(\NQ)}, because \m{\rQ} is multiplicative and the product over~$j$ of the sums~\m{\Nm_j}
is the sum over the product group;
\item \m{\Xi_{\bd}=0} unless \m{\bd\le\one};
\item for \m{j\in J} and \m{y,y'\in S}:
\[\text{if}\quad\sigma(F_j)\,y=\sigma(\Nm_j)\,y'\quad\text{and}\quad\sigma(\Nm_j)\,y=0,\qquad\text{then}\quad
\Xi_{\bd}\cdot y=\Xi_{\bd+\one_j}\cdot y'.\]
\end{enumerate}
For~(d), write \m{\Xi_{\bd}=\Xi'\cdot\sigma(S_j(\bd))}, where $\Xi'$ is the product of the other factors, the same
for $\bd$ and \m{\bd+\one_j}. If \m{\bd_j=0} then
\m{\Xi_{\bd}\,y=\Xi'\sigma(F_j)y=\Xi'\sigma(\Nm_j)y'=\Xi_{\bd+\one_j}y'}; if \m{\bd_j=1} then
\m{\Xi_{\bd}\,y=\Xi'\sigma(\Nm_j)y=0=\Xi_{\bd+\one_j}y'}; if \m{\bd_j\ge2} both sides vanish.

\emph{Step 8: the lifts do not move.} We prove by induction on the integer \m{i\ge0}: for every \m{\bd\in\NN^J}
with \m{|\bd|+i\ge|J|+1} and every \m{\Upsilon\in\EQ} with \m{\sigma(\NQ)\,\sigma(\dg{T_0})\,[\Upsilon]=0},
\begin{equation}\label{e:nomove}
\sigma\bigl(\dg{(C^2\,T_0)^{i}}\bigr)\,\Xi_{\bd}\,[\Upsilon]=0.
\end{equation}
For \m{i=0}, \m{|\bd|\ge|J|+1>|\one|}, so $\bd$ is not \m{\le\one} and \m{\Xi_{\bd}=0} by~(c). Suppose the assertion
holds for~$i$, and let $\bd$ and~$\Upsilon$ be as above with \m{i+1} in the place of~$i$. Put
\[\Pi_i=\sigma\bigl(\dg{(C^2\,T_0)^{i}}\bigr)\cdot\Xi_{\bd}\cdot C\ \in S.\]
We claim \m{\Pi_i\cdot\sigma(\rQ(\gamma))=\Pi_i} for every \m{\gamma\in\prod_j\GL_n(\Fw{w_j})}. The elements with
this property form a subgroup, so it suffices to take \m{\gamma=\gamma_j^{(j)}} with \m{j\in J} and
\m{\gamma_j\in\GL_n(\Fw{w_j})}, and to show \m{\Pi_i\cdot\sigma(\rho_j(\gamma_j)-1)=0}. Let $Z$ be as in
Step~6(i) for~$\gamma_j$. By~(d), with \m{y=C\cdot\sigma(\rho_j(\gamma_j)-1)} and \m{y'=\sigma(Z)} (the condition
\m{\sigma(\Nm_j)y=0} holds because \m{\Nm_j\rho_j(\gamma_j)=\Nm_j}),
\[\Pi_i\cdot\sigma(\rho_j(\gamma_j)-1)=\sigma\bigl(\dg{(C^2\,T_0)^{i}}\bigr)\cdot\Xi_{\bd+\one_j}\cdot\sigma(Z).\]
For every \m{\Upsilon'\in\EQ} the operator \m{Z\Upsilon'} satisfies, by Step~6(i),
\[\sigma(\NQ)\,\sigma(\dg{T_0})\,[Z\Upsilon']=[\dg{T_0}\,\NQ\,Z\,\Upsilon']=0,\]
and \m{|\bd+\one_j|+i\ge|J|+1}. So the induction hypothesis gives
\[\sigma\bigl(\dg{(C^2\,T_0)^{i}}\bigr)\,\Xi_{\bd+\one_j}\,[Z\Upsilon']=0\qquad\text{for every }\Upsilon'\in\EQ,\]
which proves the claim. By the claim and the property of~$\Psi$,
\[\Pi_i\cdot C=\sum_{\gamma}\Pi_i\,\sigma(\rQ(\gamma^{-1}))\,\sigma(\Psi)\,\sigma(\rQ(\gamma))
=\Pi_i\cdot\sigma(\Psi)\cdot\Xi_{\one},\]
using~(b). Therefore, by~(a),
\[\sigma\bigl(\dg{(C^2\,T_0)^{i+1}}\bigr)\,\Xi_{\bd}\,[\Upsilon]=\Pi_i\cdot C\cdot\sigma(\dg{T_0})\,[\Upsilon]
=\Pi_i\,\sigma(\Psi)\,\sigma(\NQ)\,\sigma(\dg{T_0})\,[\Upsilon]=0.\]

\emph{Step 9: the family kills the ideal.} Let \m{\Lambda} be the additive map from \m{\TQ[X_j:j\in J]} to~$S$ with
\[\Lambda\Bigl(\sum_{\bd}T_{\bd}\,X^{\bd}\Bigr)=\sum_{\bd\le\one}\Xi_{\bd}\cdot\sigma(\dg{T_{\bd}}),\qquad\text{so that}
\quad\Lambda(T\cdot X^{\bd})=\Xi_{\bd}\cdot\sigma(\dg{T})\ \text{ for all }T\text{ and }\bd\]
by~(c). We claim \m{2\,n!\cdot\Lambda} vanishes on~$\tie$. The ideal~$\tie$ is spanned additively by the products
of a monomial \m{T\cdot X^{\bd}} with a generator. For a generator \m{\ell^\nu\cdot\traceElem{\rr}{1}{1}{\gamma}} the
value of~$\Lambda$ is a multiple of~$\ell^\nu$, which vanishes in~$S$. For a generator \m{T_1-X_j\,T_2}, in the
notation of Step~6(ii),
\[\Lambda\bigl(T\,X^{\bd}\,(T_1-X_jT_2)\bigr)=\Xi_{\bd}\,\sigma(\dg{T_1})\,\sigma(\dg{T})
-\Xi_{\bd+\one_j}\,\sigma(\dg{T_2})\,\sigma(\dg{T}),\]
and \m{2\,n!} times this vanishes by~(d), applied to the elements
\[y=2\,n!\cdot\sigma(\dg{T_1})\,\sigma(\dg{T})\qquad\text{and}\qquad y'=2\,n!\cdot\sigma(\dg{T_2})\,\sigma(\dg{T}),\]
whose two conditions are the two assertions of Step~6(ii).

\emph{Step 10: conclusion.} By (\ref{e:ribet}), \m{T'\cdot\beta\cdot\Phi} lies in~$\tie$. In~(\ref{e:nu}) the
product \m{\upsilon\cdot2\,n!\cdot(p_0C^2)^{|J|+1}} is~$\ell^\nu$, which vanishes in~$S$. So by Step~9
and~(\ref{e:power}),
\[0=\upsilon\cdot2\,n!\cdot\Lambda(T'\beta\,\Phi)\,[\Psi]
=\upsilon\cdot2\,n!\cdot\sum_{\bd\le\one}\sigma\bigl(\dg{(C^2\,T_0)^{|J|+1}}\bigr)\,\Xi_{\bd}\,\sigma(\dg{\Phi_{\bd}})\,[\Psi].\]
Let \m{\bd<\one}. By (\ref{e:ribet}) the element \m{\Phi_{\bd}} vanishes on~$\cX$, so \m{\dg{\Phi_{\bd}}\varphi=0} by
hypothesis~(i), and (\ref{e:crit}) for \m{T''=T_0\Phi_{\bd}}, with \m{\Upsilon=0}, gives
\m{\sigma(\NQ)\,\sigma(\dg{T_0})\,[\dg{\Phi_{\bd}}\Psi]=0}. By Step~8 with \m{i=|J|+1} the term of~$\bd$ in the
sum vanishes. Only \m{\bd=\one} remains, and by (a) and~(b)
\[\sigma(\dg{T''})\,\sigma(\NQ)\,[\Psi]=0\qquad\text{for}\quad
T''=\upsilon\cdot2\,n!\cdot(C^2\,T_0)^{|J|+1}\cdot\Phi_{\one},\]
a multiple of~$T_0$. By (\ref{e:crit}) there is \m{\Upsilon\in\EQ} with
\m{\ell^\nu\cdot\Upsilon(\varphi_0)=C\cdot\dg{T''}\varphi}. By (\ref{e:power}) and~(\ref{e:nu}),
\m{T''=\ell^\nu\cdot s+\upsilon\cdot2\,n!\cdot\beta\,T'\,\Phi_{\one}}, and the second term vanishes on~$\cX$, so
that its diagonal operator kills~$\varphi$ by hypothesis~(i). Hence
\m{\ell^\nu\cdot\Upsilon(\varphi_0)=\ell^\nu\cdot C\cdot\dg{s}\varphi}, and \m{\Upsilon(\varphi_0)=C\cdot\dg{s}\varphi}.
Since $\Upsilon$ maps integral forms to integral forms and \m{\Delta\varphi_0} is integral,
\m{\Delta\cdot C\cdot\dg{s}\varphi=\Upsilon(\Delta\varphi_0)} is integral.
\end{proof}

\sect{Where the results are used}{s:use}
Outside this unit only two of its declarations are named: Theorem~A and Lemma~\ref{l:charpoly}. Both are applied
once, in the unit \emph{Selmer}, in its bound for the denominator of the projection of integral forms to lifts.
There the homomorphism~$\omega$ is the inverse of a character with values in~$\Zl$ that is trivial on the tame
generators, $\rr$ is the family of the normalised Galois representations, the places~$Q$ come from a triangular
arrangement of Taylor--Wiles primes, hypothesis~(1) comes from the Hecke polynomials of the weak lifts at the
Frobenius elements of these primes, and hypothesis~(3) is the output of Ribet's method. This is what the text of
that file shows; the note on the unit \emph{Selmer} states it. The other eight declarations serve Theorem~A only.

\sect{Notes on the formal text}{s:notes}
\subhead{Hypotheses that the census does not list.} The formal Theorem~A takes the two hypotheses of
Setting~\ref{set:aux}: the file declares them as section hypotheses and includes them, in line~87, for the
declarations that follow, and the formal proof uses both, through Proposition~\ref{p:tie}. The census reads the
arguments from the text and attaches that line to the next declaration only, which is Proposition~\ref{p:tie},
because two modifier lines stand between the inclusion and that declaration. The statement above therefore has two
hypotheses more than the census records. Lean's own answer was not obtained: nothing is compiled in this
workspace. In the same way the census does not list, for Lemma~\ref{l:test}, the natural number~$C$ on which the
base datum depends, because the lemma binds another variable of the same Lean name, the homomorphism~$\chi$.

\subhead{A statement that the census cuts.} The formal statement of Proposition~\ref{p:tie} binds \m{\rho_j} and
\m{\Nm_j} inside its type before the existential quantifier. The census ends a statement at the first such
binding, so its range (lines 96 to~99) stops before the conclusion; the statement was read to its end in line~103.

\subhead{Where the formal statements are stronger or weaker than the ones above.}
\begin{itemize}[leftmargin=1.4em,itemsep=2pt]
\item In Proposition~\ref{p:tie}, the first assertion of~(ii) holds before~$\sigma$ is applied, as the proof shows;
the formal statement asserts it after~$\sigma$, as above. In~(i) the vanishing \m{\dg{T}\cdot\NQ\cdot Z=0} is asserted in
\m{\EQ} itself, in the formal text and above.
\item In Theorem~A the function~$s$ that the proof constructs is an element of the Hecke algebra \m{\TQ}; the
formal statement, like the one above, only asserts a function. Hypothesis~(3) quantifies over all pairs
\m{(\cS,\cJ)}; the proof uses it for one pair.
\item Lemma~\ref{l:test} is stated above in Setting~\ref{set:forms} only. The formal lemma indeed does not use
that $V$ is standard, the Weil data or \m{n\ge1}: the file removes these three from it. Its last hypothesis, the
relation between $\tilde f$ and~$f$, has a default proof in the formal text and is supplied by automation at both
uses.
\item Lemma~\ref{l:charpoly} and Definition~\ref{d:trace} need neither the auxiliary places~$Q$ nor the number~$\nu$
of Setting~\ref{set:forms}; they are stated for an arbitrary family~$\bw$ indexed by~$J$, an implicit argument of
the formal text.
\item Proposition~\ref{p:first} and Lemmas \ref{l:shift} and~\ref{l:modeq} hold for every natural number~$M$,
also for \m{M=0}. The file declares \m{M\ne0} as a standing instance and removes it from each of these lemmas.
Whether the formal definition of the weighted sum carries this instance cannot be read off the text, since its
term does not use it; in every use here \m{M=\ell^\nu\ge1}.
\item In Lemma~\ref{l:line} the functions $\lgf$ and~$\delta$ are defined on all of~$\Ew{w}$ through a choice;
only their values at units of~$\Ow{w}$ matter. Hypothesis~(H5) of Proposition~\ref{p:first} is stated in the formal
text as regularity of the natural number \m{|\Gamma|} on~$R$.
\item The string length~$L$ is an implicit argument of Proposition~\ref{p:tie} and of Theorem~A, with the
hypothesis \m{L\ne0}.
\end{itemize}

\subhead{What is not used.} Of Input~8(a) only the first assertion is used, and only for the function~$0$. The
first part of the chain lemma, a lower bound for the absolute value of a root at the end of a string, is not used
in this unit. In Lemma~\ref{l:line} every assertion is used.

\subhead{The valuation.} The function \m{\ordp} of Definition~\ref{d:ord} is truncated: it is $0$ at every element
of absolute value at least~$1$. The second reader pointed this out against a reading that normalised it
at~$\ell$ alone; in Theorem~A it is applied to \m{\mu\eta} and \m{\mu\,s(\pi)}, where $\mu$ is an unrestricted
element of~$\Ql$, and to an upper bound \m{A\ge0}, with which the truncation is compatible.

\subhead{Two corrections of the second reader.} The second reader compared the dossier of this unit with the
Lean text and corrected two points, both followed here. In the proof of Lemma~\ref{l:test} the form
\m{\varphi[\psi]} vanishes whenever \m{\psi^{-1}(e)=\theta}, and it may vanish otherwise; the hypothesis is applied
in the non-zero case, where \m{\psi^{-1}(e)\ne\theta} follows. And the valuation is the truncated one. The second
reader also asked that the tie keep both the sum over the group and the factor \m{2\cdot n!}, and that the two
different trace elements be kept apart; Proposition~\ref{p:tie} does so.

\subhead{What the authors' comments name as related work.} The head of the second file says that auxiliary tame
level and first-order derivatives in diamond operators occur in the literature in other settings and names
A.~Wiles (1990), E.~Lecouturier (2019) and A.~Dotto and J.~Thorne (2026); for the chain lemma it names M.~Tadić
(1986), W.~Casselman (1980) and L.~Clozel, M.~Harris and R.~Taylor (2008, Lemma~3.1.3 and §3.4). These attributions
are the authors'; the references at the end repeat them as the file gives them, and the catalogue of the handout
says which places were confirmed.

\subhead{How this note was written.} From the unit's dossier: the reader's statements of the \UnitMembers{}
declarations and of the imports, and the six structured proofs, every step; and from the second reader's message on
this unit. The two Lean files of the unit were read in full, statements and proofs, and each statement above was
compared with its Lean text argument by argument. In the library the following were read in the Lean text, as
statements and definitions only: the base datum, the levels, the integral forms, the eigensystems and the
Taylor--Wiles places; the diagonal operators and the Hecke algebra; the operators at the auxiliary places with
Inputs~2(c) to~(f); the root strings and the chain lemma; the two statements of Input~8 to their ends, which the
census cuts; the two definitions of Definitions \ref{d:ideal} and~\ref{d:ribet}; the valuation; the two statements
of Input~7; the filter of Input~3. No proof of the library was opened. The reader marked three places as least
certain. They were derived again from the Lean lines: the linear form in the proof of hypothesis~(H3) of the
first-order lemma inside Proposition~\ref{p:tie} (lines 132 to~151: the form \m{\varphi'\mapsto\ip{\dg{T}\varphi'}{\varphi_1}},
tested on base-level eigenforms, as the reader says); the induction of Step~8 (lines 264 to~279); and Steps 9
and~10 (lines 280 to~304), including the cancellation of the factors \m{(p_0C^2)^{|J|+1}} and~$\upsilon$. The
reader's claims agree with the Lean text in all three. The arguments above follow the formal ones, with these
differences: the transfer of a central element is computed from formula~(\ref{e:transfer}) and not quoted; the
module of Step~5 is presented as a quotient by a left ideal, which is what the formal construction is; and the
product map through which the formal text defines the operators \m{\Xi_{\bd}} has no name here. The program
check of the unit's dossier covered \UnitSteps{} leaf steps of \UnitProofs{} proofs. Nothing was compiled for this
note.

\head{Appendix A. Concordance with the Lean artifact}
{\fontsize{8.6}{11.6}\selectfont\renewcommand{\leanfont}{\ttfamily\fontsize{7.7}{10}\selectfont\color{leancol}}
\setlength{\LTleft}{0pt}\setlength{\LTright}{\fill}\setlength{\tabcolsep}{0pt}
\begin{longtable}{@{}p{30mm}p{56mm}p{43mm}r@{}}
\textbf{In this note} & \textbf{Lean declaration} & \textbf{File} & \textbf{Lines}\\\midrule\endhead
\ConcordanceRows
\end{longtable}}

\subhead{Theorem A and the first-order lemma as they are written in the artifact.}
\quoted{mainThm}
\quoted{firstOrder}
In these texts {\leanfont EllAdicIntegers} is~$\Zl$, {\leanfont PadicAlgCl ℓ} is~$\Ql$ and {\leanfont Ω}
is~$\Qbar$; {\leanfont D} is the base datum~$\fD$ and {\leanfont c} the number~$C$; {\leanfont Q.w} is the
family~$\bw_Q$ and {\leanfont ![]} the empty family; {\leanfont D.Eigensystem}, {\leanfont D.heckeAlgebra},
{\leanfont D.integralEnd}, {\leanfont D.formsOfLevel}, {\leanfont D.localOps} and {\leanfont D.residueRep} are
\m{\Sys_{\fD}}, \m{\Hk_{\fD}}, \m{\Eint_{\fD}}, \m{\Forms_{\fD}}, \m{\Loc_{\fD}} and \m{\rQ}; {\leanfont diagOp} and
{\leanfont heckeAction} are both written \m{\dg{\,\cdot\,}}; {\leanfont ellVal} is \m{\ordp}; {\leanfont card n} is the
rank~$n$; {\leanfont N} is the family~$\rr$ and {\leanfont a} the homomorphism~$\omega$; {\leanfont b}, {\leanfont e},
{\leanfont t}, {\leanfont m} and {\leanfont γ} of Theorem~A are $\beta$, $\varepsilon$, $L$, $\mu$ and~$A$;
{\leanfont l}, {\leanfont d} and {\leanfont y} in its conclusion are $\varphi_0$, $\Delta$ and~$\varphi$;
{\leanfont HasRootString}, {\leanfont IsTraceIdeal} and {\leanfont IsRibetElement} are the notions of Definitions
\ref{d:string}, \ref{d:ideal} and~\ref{d:ribet}. In the first-order lemma {\leanfont a}, {\leanfont ac},
{\leanfont L}, {\leanfont γ}, {\leanfont n}, {\leanfont k} and {\leanfont r} are $\alpha$, the family \m{(c_i)},
$\Pi$, $\gamma_0$, $C$, $e$ and~$y$, and {\leanfont a.transfer} is \m{\Tr_\alpha}. In the formal
Proposition~\ref{p:tie} the ring called {\leanfont R} is the ring~$S$ above.

\head{Notation}
\begingroup\fontsize{8.6}{10.9}\selectfont\raggedright
\setlength{\columnsep}{9mm}\setlength{\parskip}{0pt}
\begin{multicols}{2}
\nt{$\ell$, $\iota$, $\Ql$, $\Zl$}{a prime number; the fixed isomorphism \m{\CC\to\Ql}; an algebraic closure
of~$\mathbb{Q}_\ell$; its ring of integers}
\nt{$|z|$, $\ordp(z)$}{the absolute value of~$\Ql$ with \m{|\ell|=\ell^{-1}}; the truncated valuation
(Definition~\ref{d:ord})}
\nt{$E$, $\Qbar$, $\GE$}{a number field, from Setting~\ref{set:forms} on a CM field; an algebraic closure
of~$\mathbb{Q}$; the Galois group of~$\Qbar$ over~$E$}
\nt{$w$, $\bar w$, $\Ew{w}$, $\Ow{w}$, $\pw{w}$, $\Fw{w}$, $q_w$}{a finite place of~$E$; its conjugate; the
completion; its integers; their maximal ideal; the residue field; its number of elements}
\nt{$n$, $1_n$}{the rank, the size of the matrices; the identity matrix}
\nt{$V$, $\kappa$, $\Gf$, $\AllForms$}{the hermitian space; the weight; the finite adelic points of its unitary
group; the forms}
\nt{$\loc_w$, $\Par_w$, $\Iw_w$}{the representation of \m{\GL_n(\Ew{w})} on the forms; two subgroups of
\m{\GL_n(\Ow{w})} (Definition~\ref{d:par})}
\nt{$\varpi$, $u$, $U_\gamma$, $U_\varpi$}{a uniformiser; a unit of~$\Ow{w}$; Hecke operators of double cosets of
\m{\Iw_w}}
\nt{$C$, $\fD$}{a natural number; a base datum for \m{(V,\kappa,C)}; in Proposition~\ref{p:first}, $C$ is any
natural number}
\nt{$\bw$, $\emptyset$, $\Forms_{\fD}(\bw)$, $\Eint_{\fD}(\bw)$}{a family of finite places; the empty family; the
forms of level~$\bw$; their integral endomorphisms}
\nt{$\xi$, $\pi$, $\Sys_{\fD}(\bw)$, $\cX$}{isotypic spaces, eigensystems; the eigensystems of level~$\bw$; a set
of eigensystems}
\nt{$\hp{\xi}{w}$, $\theta_{\xi,w,i}$}{the Hecke polynomial of~$\xi$ at~$w$; the chosen Hecke eigenvalues}
\nt{$\theta$, $\theta_i$, $\theta'_i$}{a complex number, or a function with complex values; in Input~8,
characters}
\nt{$s$, $s'$, $\dg{s}$}{functions from eigensystems to~$\Ql$; the diagonal operator}
\nt{$T$, $T'$, $T''$, $T_\xi$, $T_{*}$, $T_0$, $T_1$, $T_2$, $T_G$, $T^{\tau}_G$, $T_{\bd}$}{elements of a Hecke
algebra \m{\Hk_{\fD}(\bw)}}
\nt{$J$, $j$, $j'$, $\nu$, $Q$, $w_j$, $q_j$, $\bw_Q$}{a finite set; its elements; an exponent of~$\ell$; the
Taylor--Wiles places; the place of index~$j$; its norm; the family of the places}
\nt{$\Fr_j$, $\tau_j$, $\Art_w$}{the geometric Frobenius element and the tame generator at~$w_j$; the local Artin
map}
\nt{$\rQ$, $\NQ$, $\gamma^{(j)}$, $\LQ{j}$}{the action of \m{\prod_j\GL_n(\Fw{w_j})}; the sum over this group;
the element with $\gamma$ at~$j$; the operators at~$w_j$}
\nt{$\rho_j$, $\Nm_j$}{the action of \m{\GL_n(\Fw{w_j})}; the sum over this group}
\nt{$\varrho_\xi$, $\omega$, $\rr$, $\rr_\xi$, $\rr_i$}{the chosen Galois representation; a homomorphism
\m{\GE\to\Zl}; the twisted family; its members}
\nt{$\traceElem{\rr}{G}{\gamma'}{\gamma}$}{the trace element (Definition~\ref{d:trace})}
\nt{$G$, $G'$, $G''$, $G_1$, $G^{\iota}$, $X$, $X_j$}{polynomials in~$X$; the variable; the variables of Ribet's
element}
\nt{$\Gamma$, $\gamma$, $\gamma'$, $\gamma''$, $\gamma_0$, $\gamma_j$, $\bar\gamma$, $\gamma_{11}$}{a finite group;
elements of a group; the distinguished element of Lemma~\ref{l:line}; a reduction; a first entry}
\nt{$H$, $h$, $h'$, $h_{11}$, $\alpha$, $\Tr_\alpha$}{a subgroup; its elements; a first entry; a homomorphism
\m{H\to\ZZ/M}; its transfer}
\nt{$M$, $z$, $z'$, $z_i$, $\bar z$}{a natural number, the modulus; arguments of functions: residues, elements of
a field; a reduction}
\nt{$\lgf$, $\lgf_0$, $\delta$, $\delta_0$}{the functions of Lemma~\ref{l:line} and the homomorphisms that define
them}
\nt{$R$, $S$, $\sigma$, $1_S$}{rings; a ring homomorphism \m{R\to S}; the unit of~$S$}
\nt{$\rho$, $\chi$, $\psi$}{a homomorphism from a group to a ring; a homomorphism to \m{\ZZ/M}; an additive
character}
\nt{$f$, $f'$, $f''$, $f_0$, $\tilde f$}{functions with integral or complex values}
\nt{$\weightedSum{\rho}{\chi}{f}$, $\dif{e}f$}{the weighted sum (Definition~\ref{d:wsum}); the difference
\m{f(\,\cdot-e)-f}}
\nt{$B$, $\cA$, $I$, $i$, $i'$, $c_i$}{an additive map on functions; a subring; an index set; indices, also of
rows and columns; a commuting family}
\nt{$\Psi$, $\Pi$, $\Pi_i$, $F$, $F_j$, $Z$, $y$, $y'$}{elements of a ring: a Higman operator; an invariant
element; the operator of the tie; elements}
\nt{$e$, $L$}{a natural number or its class: a step, an exponent; the length of a root string}
\nt{$\Sigma$, $\varsigma$, $\rep{\gamma}$, $\Hpart{\gamma}$, $\Yp$, $\Ym$, $\Nm$, $v$}{representatives of cosets;
one of them; the representative and the $H$-part of~$\gamma$; the two coset sums; the sum over~$\Gamma$; the least
absolute residue}
\nt{$Y$, $Y'$, $\cB$}{operators; a linear map (Lemma~\ref{l:test})}
\nt{$\varphi$, $\varphi'$, $\varphi_0,\dots,\varphi_3$, $\varphi[f']$, $\ip{\varphi'}{\varphi}$}{forms; the
invariant hermitian form}
\nt{$\beta$, $p$, $p_0$, $\varepsilon$, $\varepsilon_j$, $\varepsilon^{+}$, $\varepsilon^{\sharp}$}{a function on
eigensystems; a polynomial; its constant coefficient; polynomials over~$\Zl$}
\nt{$\mu$, $A$, $\eta$, $\Delta$, $\upsilon$}{an element of~$\Ql$; a bound in \m{[0,\infty]}; a Ribet coefficient;
a natural number; an element of~$\Zl$}
\nt{$\cS$, $\cJ$, $\Phi$, $\Phi_{\bd}$, $\tie$, $\ev$}{a subalgebra; a submodule; Ribet's polynomial; its
coefficients; the ideal of the ties; evaluation}
\nt{$\bd$, $\bd_j$, $|\bd|$, $\one$, $\one_j$, $X^{\bd}$}{a multi-index; its components; their sum; the
multi-indices of \m{\prod_jX_j} and of~$X_j$; a monomial}
\nt{$\nn$, $\cM$, $[\Upsilon]$, $\Upsilon$, $\Upsilon'$}{a left ideal; the quotient module; a class; integral
operators}
\nt{$\Xi_{\bd}$, $\Xi'$, $S_j(\bd)$, $\Lambda$}{the operators of Step~7; an additive map (Step~9)}
\end{multicols}
\endgroup

\head{References}
{\fontsize{9.4}{12.6}\selectfont
The places are those that the catalogue of the handout and the second reader's comparison give; where they say
that a place is not confirmed, so does the text above. No reference was typed from memory.
\begin{itemize}[leftmargin=1.4em,itemsep=1pt]
\item J.~Bernstein, Lecture Notes in Mathematics 1041 (1984).
\item G.~Chenevier, M.~Harris, Construction of automorphic Galois representations~II, Cambridge J.\ Math.\ 1
(2013), Theorem~3.2.3.
\item L.~Clozel, M.~Harris, R.~Taylor, Automorphy for some l-adic lifts of automorphic mod~l Galois
representations, Publ.\ Math.\ IHÉS 108 (2008); §3.3, Lemma~3.3.1 (second reader, for Input~2(f)); Lemma~3.1.3 and
§3.4 (the authors, for the chain lemma).
\item M.~Harris, R.~Taylor, On the geometry and cohomology of some simple Shimura varieties (2001).
\item D.~G.~Higman, Modules with a group of operators, Duke Math.\ J.\ 21 (1954).
\item J.-P.~Serre, Local Fields (1979), Chapters XII to~XIV; the second reader gives XIII~§4, Proposition~13 with
its corollary, and XIV~§6, Theorem~1 and Corollary~2.
\item M.~Tadić, Ann.\ Sci.\ École Norm.\ Sup.\ 19 (1986).
\item R.~Taylor, A.~Wiles, Ring-theoretic properties of certain Hecke algebras, Ann.\ of Math.\ 141 (1995).
\item R.~Taylor, T.~Yoshida, Compatibility of local and global Langlands correspondences, J.\ Amer.\ Math.\
Soc.\ 20 (2007), Theorem~A.
\item Named by the authors' comments only, as the file gives them: A.~Wiles, The Iwasawa conjecture for totally
real fields (1990); E.~Lecouturier, Mazur--Tate conjecture for prime conductor (2019); A.~Dotto, J.~Thorne,
Residually reducible representations of rank~2 (2026); W.~Casselman, Unramified principal series of $p$-adic
groups~I (1980).
\end{itemize}}

\end{document}
